\documentclass[11pt,reqno]{amsart}
\usepackage[a4paper,margin=27mm]{geometry}
\usepackage[T1]{fontenc}
\usepackage{lmodern,microtype}
\usepackage{amsmath,amssymb,mathtools}
\usepackage{booktabs,array,longtable}
\usepackage{enumitem,needspace}
\usepackage[hidelinks]{hyperref}
\usepackage{aliascnt}
\newtheorem{theorem}{Theorem}[section]
\newaliascnt{proposition}{theorem}
\newtheorem{proposition}[proposition]{Proposition}
\aliascntresetthe{proposition}
\newaliascnt{lemma}{theorem}
\newtheorem{lemma}[lemma]{Lemma}
\aliascntresetthe{lemma}
\newaliascnt{corollary}{theorem}
\newtheorem{corollary}[corollary]{Corollary}
\aliascntresetthe{corollary}
\theoremstyle{definition}
\newaliascnt{definition}{theorem}
\newtheorem{definition}[definition]{Definition}
\aliascntresetthe{definition}
\theoremstyle{remark}
\newaliascnt{remark}{theorem}
\newtheorem{remark}[remark]{Remark}
\aliascntresetthe{remark}
\usepackage[nameinlink,noabbrev]{cleveref}
\crefname{theorem}{Theorem}{Theorems}
\crefname{proposition}{Proposition}{Propositions}
\crefname{lemma}{Lemma}{Lemmas}
\crefname{corollary}{Corollary}{Corollaries}
\crefname{definition}{Definition}{Definitions}
\crefname{remark}{Remark}{Remarks}
\DeclareMathOperator{\Exc}{Exc}
\DeclareMathOperator{\Sing}{Sing}
\DeclareMathOperator{\rk}{rk}
\DeclareMathOperator{\Pic}{Pic}
\DeclareMathOperator{\Supp}{Supp}
\DeclareMathOperator{\supp}{supp}
\DeclareMathOperator{\Null}{Null}
\DeclareMathOperator{\NS}{NS}
\DeclareMathOperator{\Irr}{Irr}
\DeclareMathOperator{\wt}{wt}
\DeclareMathOperator{\ord}{ord}
\DeclareMathOperator{\mult}{mult}
\DeclareMathOperator{\Bl}{Bl}
\DeclareMathOperator{\tr}{tr}
\newcommand{\eps}{\varepsilon}
\newcommand{\PP}{\mathbb P}
\newcommand{\FF}{\mathbb F}
\newcommand{\QQ}{\mathbb Q}
\newcommand{\ZZ}{\mathbb Z}
\newcommand{\RR}{\mathbb R}
\newcommand{\cO}{\mathcal O}
\newcommand{\ADE}{\mathrm{ADE}}
\numberwithin{equation}{section}
\allowdisplaybreaks[2]
\setcounter{tocdepth}{1}
\hypersetup{
 pdftitle={Singular points of klt del Pezzo surfaces of Picard number one in positive characteristic},
 pdfsubject={A seven-point bound in odd positive characteristic and Frobenius constructions},
 pdfkeywords={del Pezzo surfaces, Kawamata log terminal singularities, Picard number, positive characteristic, rational surfaces}}
\title[Singular points of klt del Pezzo surfaces]{Singular points of klt del Pezzo surfaces of Picard number one in positive characteristic}
\date{}
\begin{document}
\begin{abstract}
Let $X$ be a normal projective surface over an algebraically closed field
of odd positive characteristic. Suppose that $X$ has Kawamata log
terminal singularities, that $-K_X$ is ample and $\QQ$-Cartier, and that
its Picard number is one. We prove that $X$ has at most seven distinct
singular points. A construction by Bernasconi gives such a surface with
seven singular points in characteristic three, so the bound is sharp
in that characteristic. The proof proceeds on the minimal resolution
of a hypothetical counterexample. Among its smooth rational curves of
self-intersection $-1$, we choose one of least degree with respect to
the pullback of $-K_X$. Effective divisors obtained from Riemann--Roch,
together with rational fibrations, reduce the possible configurations
of curves contracted by the resolution to a finite list. Intersection theory and
branched double covers exclude each configuration. We also describe a
construction using the Frobenius morphism that recovers Bernasconi's
surface and the characteristic-two families of Keel and McKernan.
These characteristic-two constructions give surfaces with Kawamata log
terminal singularities, ample anticanonical divisor, Picard number one,
and arbitrarily many singular points.
\end{abstract}
\maketitle

\section{Introduction}\label{sec:introduction}
A \emph{klt del Pezzo surface} over an algebraically closed field $k$ is a
normal projective surface $X$ with Kawamata log terminal singularities
such that $-K_X$ is an ample $\QQ$-Cartier divisor. Here Kawamata log
terminal, abbreviated klt, means that all discrepancies of the pair
$(X,0)$ are greater than $-1$. The Picard number $\rho(X)$ is the dimension
of the real vector space of Cartier divisor classes modulo numerical
equivalence. We use \emph{rank one} only to mean Picard number one. The assumption
$\rho(X)=1$ restricts the numerical geometry
of $X$, but does not directly bound the number of irreducible curves
appearing on a resolution of its singularities. We consider the number
\[
 n(X):=\#\Sing(X)
\]
of distinct singular points, rather than the number of exceptional
curves or the length of a singular subscheme.

\begin{theorem}[Uniform seven-point bound]\label{thm:main}
Let $k$ be an algebraically closed field of characteristic $p>2$, and let
$X$ be a normal projective surface over $k$. Suppose that $X$ has klt
singularities, $-K_X$ is ample and $\QQ$-Cartier, and $\rho(X)=1$.
Then $n(X)\le7$.
\end{theorem}

There is no lower bound on the discrepancies in this statement other
than the strict inequalities in the definition of klt. In particular,
no fixed Cartier index is assumed. The characteristic-three case of
\cref{thm:main} is sharp. In characteristic two the number of singular
points is unbounded. The existence statements below are due to Bernasconi
\cite[Section~3.1 and Proposition~3.2]{Bernasconi} and Keel--McKernan
\cite[Section~9]{KeelMcKernan}, respectively.

\begin{theorem}[Bernasconi; Keel--McKernan]\label{thm:examples}
Over every algebraically closed field of characteristic three, there is
a klt del Pezzo surface $X$ with $\rho(X)=1$, $n(X)=7$, and $K_X^2=1/3$.
Over every algebraically closed field of characteristic two, for each
integer $n\ge3$ there is a klt del Pezzo surface $X_{2,n}$ with
\[
 \rho(X_{2,n})=1,\qquad n(X_{2,n})=2n+1,\qquad
 K_{X_{2,n}}^2=\frac{2}{n-2}.
\]
\end{theorem}

\subsection{Related work}\label{subsec:related-work}
Over $\mathbb C$, Keel and McKernan obtained
singular-point estimates from their study of rational curves and
orbifold inequalities \cite[Section~9]{KeelMcKernan}. Belousov proved
that a klt del Pezzo surface of Picard number one has at most four
singular points \cite[Theorem~1.1]{Belousov}. Liu and Xie subsequently
established the inequality
\[
 n(X)\le 2\rho(X)+2
\]
for complex klt Fano surfaces of arbitrary Picard number
\cite[Theorem~1.2]{LiuXie}. These results concern characteristic zero.

In positive characteristic, Lacini's classification of rank-one log del
Pezzo surfaces in characteristics different from two and three gives
$n(X)\le4$ for $p>5$ \cite[Corollary~1.2]{Lacini}. Lacini also constructs a surface with five singular points in
characteristic five \cite[Theorem~1.3 and Example~7.6]{Lacini}. Nagaoka's classification of the
characteristic-five surfaces whose resolutions with exceptional
boundary do not lift to the ring of Witt vectors gives a more detailed
description of the exceptional behavior in that characteristic
\cite{NagaokaFive}. His later review re-examines the classification
arguments and their characteristic restrictions, and identifies a
possible omission from one of Lacini's lists \cite{NagaokaReview}.
The proof below does not use these classifications. For $p>5$, the known four-point bound is stronger than seven.

Keel and McKernan constructed rank-one klt del Pezzo surfaces in
characteristic two with $2n+1$ singular points for arbitrarily large
$n$ \cite[Section~9]{KeelMcKernan}. Bernasconi constructed a
seven-point surface in characteristic three in his study of the failure
of Kawamata--Viehweg vanishing \cite[Section~3.1]{Bernasconi}; the
same surface is recorded in \cite[Example~7.4]{Lacini}. Its minimal
exceptional divisor consists of four isolated $(-3)$-curves and three
chains of two $(-2)$-curves. The graph-of-Frobenius construction used in
\cref{sec:examples} is Bernasconi's construction when $p=3$ and $n=3$,
up to interchanging the factors of $\PP^1\times\PP^1$.
Bernasconi also noted the reversal of the canonical sign for its
higher-characteristic analogues \cite[Remark~3.3]{Bernasconi}.
\Cref{sec:examples} records a uniform canonical-divisor formula and a
contraction proof using an explicit nef divisor with its exact null
locus. The further calculations describe the shortest curves and their
birational replacements.

Palka and Pe\l ka pursue a classification over arbitrary algebraically
closed fields by means of the \emph{height}: the least intersection
number of the reduced exceptional divisor with a fiber of a rational
ruling on the minimal resolution, with value $+\infty$ if there is no
such ruling. Their first paper treats heights one
and two and surfaces admitting specified birational morphisms to
canonical del Pezzo surfaces \cite{PalkaPelkaI}; the next treats height
three \cite{PalkaPelkaIII}. In \cite[Remark~1.7]{PalkaPelkaI},
Palka and Pe\l ka already announce that the maximal numbers of singular
points in characteristics three and five are seven and five,
respectively, with proofs deferred to subsequent articles. The proof below establishes the seven-point bound in every odd positive characteristic.

The operation of removing a contacted exceptional $(-2)$-curve
and contracting a nonexceptional $(-1)$-curve, under a two-contact
hypothesis, is an elementary swap in
\cite[Definition~1.3 and Lemma~2.9]{PalkaPelkaI}. The replacements in
\cref{sec:replacements} include these operations. We require a
formulation with explicit discrepancy coefficients and exact changes
in the number of singular points, also allowing the additional contact
configurations arising at a curve of least degree.

Dolgachev, Mendes Lopes, and Pardini associate a binary code to disjoint
$(-2)$-curves by recording which sums are twice a Picard class, and use
coverings to constrain nodal rational surfaces
\cite[Sections~2--3]{DolgachevMendesLopesPardini}. In particular, the
fact that every such sum contains a multiple of four curves is the
standard Riemann--Roch parity restriction.
Calabri, Ciliberto, and Mendes Lopes show that an even set of four nodes
on a smooth rational surface in characteristic different from two
forces a rational fibration with two distinguished fibers
\cite[Theorem~1.1]{CalabriCilibertoMendesLopes}. Here the double-cover argument keeps track
of the entire exceptional forest, including the components outside the
branch divisor, and combines its rank with the Picard-number-one
hypothesis. The characteristic-two divisibility calculation identifies
the classical code $DE(n)$ of
\cite[Section~2]{DolgachevMendesLopesPardini} on the explicit Frobenius
surface.

We work on one fixed minimal resolution and choose a curve of least
degree with respect to the pullback of the anticanonical divisor. Effective integral adjoints give degree inequalities
without a vanishing theorem, and reduce the problem to a small
classification of weighted forests rather than a classification of
surfaces. The external contraction results are Keel's semiampleness
criterion \cite[Theorem~0.2]{Keel} and the smooth-surface case of
Tanaka's contraction theorem \cite[Theorem~4.4]{Tanaka}.
Within positive characteristic, the remaining restriction is precisely
the separability of two degree-two constructions. Consequently the
characteristic-three proof works unchanged for each fixed $p>2$;
there is no specialization or lifting argument from characteristic
three. 

\subsection{Outline of the proof}
Let $\pi:S\to X$ be the minimal resolution, let
$D=\Exc(\pi)_{\mathrm{red}}$, and put $L=\pi^*(-K_X)$. The standard resolution properties give a rational surface
$S$ and a simple-normal-crossing forest $D$ of smooth rational curves. Its irreducible components are independent in
$\Pic(S)$ and number $\rho(S)-1$. The connected components of this
forest correspond exactly to the singular points of $X$.

A curve on $S$ not contained in $D$ will be called \emph{exterior}.
Suppose a counterexample to \cref{thm:main} exists over a fixed field,
and choose one with $\rho(S)$ minimal. Choose an exterior $(-1)$-curve
$P$ for which $\ell=L\cdot P$ is least. An extremal-ray argument on the
smooth surface proves that $L+\ell K_S$ is nef. Consequently every
exterior integral curve $Q$ satisfies
\[
 L\cdot Q\ge\ell(-K_S\cdot Q)\ge\ell.
\]
This controls the exterior components of effective integral adjoint
divisors, not only the $(-1)$-curves initially under consideration.

We use adjoint divisors of the form $K_S+A$, where $A$ is an effective
integral divisor formed from $P$ and exceptional curves. Riemann--Roch
produces effective representatives of these adjoint classes. Their small $L$-degree either
forces a singular-point bound or produces a shorter exterior
$(-1)$-curve. When $P$ meets two disjoint exceptional $(-2)$-curves,
they form a rational fiber with $P$ of multiplicity two. A bisection is
a curve whose morphism to the base of the ruling has degree two.
The fiber intersection equations and the adjoint of a bisection then
give the required bound. Certain remaining contact configurations admit
birational replacements that do not decrease the number of singular
points and strictly lower the Picard number of the minimal resolution.
The replacements are realized by projective contractions; their
anticanonical divisors and discrepancies are computed explicitly.

These arguments leave one possibility. There are pairwise disjoint
exceptional curves $C,B_1,B_2$ of respective self-intersections
$-2,-3,-\beta$, with $\beta\in\{3,4,5\}$, and an exterior $(-1)$-curve
$R$ disjoint from $C+B_1+B_2+P$, such that
\[
 R\sim K_S+C+B_1+B_2+2P.
\]
This identity bounds the weights and the number of edges of the entire
exceptional forest. A direct classification gives ten possibilities.
The full-rank lattice generated by $D$ and $P$ must have square
absolute determinant, since the Picard lattice of a rational surface
is unimodular. Seven possibilities fail this test. Riemann--Roch
parity excludes two more; the last forces four isolated exceptional
$(-2)$-curves whose sum is twice a Picard class. A separable double
cover branched over such a sum contradicts the Hodge index theorem.

The bound is sharp in characteristic three. In every odd
positive characteristic, a degree-two map from a rational bisection is
separable and tame, and the cover branched over an even set of isolated
exceptional nodes is separable. All other characteristic-dependent
inputs used here hold in every positive characteristic. The examples
in \cref{sec:examples} exhibit the failure of both degree-two mechanisms
in characteristic two. Since the contractions use Keel's
positive-characteristic semiampleness theorem, this argument is not
asserted to extend to characteristic zero.

\tableofcontents

\section{Resolutions, intersection theory, and projective contractions}
\label{sec:preliminaries}
\subsection{Conventions}
Unless a different field is specified, $k$ is algebraically closed of
positive characteristic. A surface is integral and projective over $k$.
An integral curve is reduced and irreducible. On a smooth surface a
$(-b)$-curve means a smooth rational curve of self-intersection $-b$.
We write $\sim$, $\sim_{\QQ}$, and $\equiv$ for linear, rational linear,
and numerical equivalence, respectively. A rational divisor is nef if
it has nonnegative intersection with every integral curve, and it is
semiample if a positive integral multiple is basepoint free. A nef
divisor on a smooth surface is big precisely when its square is
positive. Its null locus is the reduced union of curves of degree zero.
We write $N^1(S)_{\QQ}$ for rational divisor classes modulo numerical
equivalence, and $\overline{\mathrm{NE}}(S)$ for the closed cone of
effective curve classes.

For a minimal resolution $\pi:S\to X$ of a klt del Pezzo surface of
Picard number one, set
\[
 D=\Exc(\pi)_{\mathrm{red}}=\sum_iD_i,\qquad
 L=\pi^*(-K_X),\qquad v=L^2=K_X^2.
\]
This triple will be called a \emph{resolution datum}. An exterior curve
is an integral curve on $S$ not contained in $D$. The canonical
pullback will be written
\[
 K_S+\sum_i\lambda_iD_i=\pi^*K_X=-L.
\]
Thus the discrepancy of $D_i$ is $-\lambda_i$, and its log discrepancy
is $1-\lambda_i$. We prove below that $0\le\lambda_i<1$.
A coefficient attached to a $(-2)$-curve need not vanish when that
curve is in a connected exceptional divisor containing a curve of
smaller self-intersection. We therefore distinguish self-intersection
from discrepancy throughout.

For a reduced simple-normal-crossing divisor $G=\sum G_i$, its dual
graph has one vertex for each component and an edge for each
intersection point. Simple normal crossings means that all components
are smooth, every intersection is transverse, and no three components
pass through one point. A \emph{forest} has no cycles; its connected
components are trees. We give its vertices weights $b_i=-G_i^2$ and
write
\[
 A_G=-(G_i\cdot G_j),\qquad q_i=K_S\cdot G_i=b_i-2
\]
when the components are rational. We call a vertex \emph{weight-two}
when $b_i=2$ and \emph{higher-weight} when $b_i\ge3$. Contact vectors
use the intersection numbers with the individual integral components,
including their local intersection multiplicities. The notation
$[b_1,\ldots,b_m]$ denotes a chain with those weights; a bracket
containing a curve name specifies that vertex rather than its weight.

We use the Hodge index theorem, adjunction, Riemann--Roch and Serre
duality on smooth projective surfaces, as well as contraction of
$(-1)$-curves and factorization of birational morphisms between smooth
surfaces into point blowups. Our use of the surface minimal model
program is over $\operatorname{Spec}k$ and with zero boundary on a
smooth surface; the needed contraction and termination results are
\cite[Theorems~4.4 and~1.1]{Tanaka}. We also use the classical
ruledness criterion in positive characteristic: a smooth projective
surface of Kodaira dimension $-\infty$ is birationally ruled, with
rational surfaces included. In particular a smooth surface with
non-pseudo-effective canonical divisor has this property: it has
$h^0(S,12K_S)=0$, so one may apply
\cite[Theorem~2.1\textup{(I)}]{CataneseLi}.
We write $p_g(S)=h^0(S,K_S)$ and $q(S)=h^1(S,\cO_S)$; here $q$ is
the arithmetic irregularity, which agrees with the genus of the base
on a birationally ruled surface. For an integral curve $C$ on a smooth
surface, $p_a(C)=1+(C^2+K_S\cdot C)/2$ denotes its arithmetic genus.

\subsection{Inverse intersection matrices and rational trees}
We recall two standard facts used in the contraction arguments.
\begin{lemma}[Positivity of a Stieltjes inverse]\label{lem:stieltjes}
Let $A$ be a symmetric positive-definite matrix with nonpositive off-diagonal entries.  Then $A^{-1}$ is entrywise nonnegative.  If the graph of the nonzero off-diagonal entries is connected, then $A^{-1}$ is entrywise positive.
\end{lemma}

The first assertion is the usual inverse-positivity property of a
positive-definite Stieltjes matrix; the strict assertion is its
irreducible case.

\begin{lemma}[Picard group of a rational tree]\label{lem:tree-picard}
Let $Z$ be a connected reduced nodal curve whose irreducible components are copies of $\PP^1$ and whose dual graph is a tree.  The multidegree map
\[
  \Pic(Z)\longrightarrow\ZZ^{\{\text{components of }Z\}}
\]
is an isomorphism.  In particular, a line bundle having degree zero on every irreducible component of $Z$ is trivial.
\end{lemma}

This is the standard normalization description of the Picard group of
a nodal curve: the gluing factor is $(k^*)^{b_1}$, where $b_1$ is the
first Betti number of the dual graph, and is trivial for a tree.

\subsection{Semiampleness and the minimal resolution}

\begin{theorem}[Keel {\cite[Theorem~0.2]{Keel}}]\label{thm:keel}
Let $M$ be a nef line bundle on a scheme projective over a field of positive characteristic.  Let $E(M)$ be the reduced union of the positive-dimensional subvarieties on which $M$ is not big.  Then $M$ is semiample if and only if $M|_{E(M)}$ is semiample.
\end{theorem}



The following collects standard resolution properties. Rationality is
\cite[Lemma~5.1]{Bernasconi}. We give the discrepancy and rank arguments
in the form needed for the subsequent contractions, without invoking
a classification of quotient singularities.

\begin{proposition}[Geometry of the minimal resolution]
\label{prop:minimal-resolution}
Let $X$ be a rank-one klt del Pezzo surface over an algebraically closed
field of positive characteristic, and let $\pi:S\to X$ be its minimal
resolution. Put $D=\Exc(\pi)_{\rm red}$ and $L=\pi^*(-K_X)$.
Then $D$ is an SNC forest of smooth rational curves of square at most
$-2$. Its negative intersection matrix $A$ is positive definite, and
\[
 K_S+\sum_i\lambda_iD_i=\pi^*K_X,\qquad
 \lambda=A^{-1}q,\quad q_i=-2-D_i^2,\quad 0\le\lambda_i<1.
\]
The surface $S$ is rational, $\Pic(S)$ is torsion-free and unimodular,
$K_S^2+\rho(S)=10$, and $\#\Irr(D)=\rho(S)-1$. Moreover
$\#\pi_0(D)=n(X)$.
In particular, $S$ admits a birational morphism to $\PP^2$ or a
rational ruling. 
\end{proposition}

\begin{proof}
The resolution is an isomorphism over the smooth locus of $X$ and
has connected fibers. Its exceptional connected components therefore
correspond exactly to the singular points of $X$.
First the exceptional classes are independent. A numerical relation,
after separating signs, would give numerically equivalent effective
exceptional divisors $Z_+,Z_-$ with no common component. If both are
nonzero then $Z_+^2=Z_+\cdot Z_-\ge0$, contradicting Hodge index since
$L\cdot Z_+=0$ and $L^2>0$. A relation of only one sign is excluded by
an ample divisor. Hodge index now makes their intersection matrix
negative definite. For an exceptional integral curve $E$ put $b=-E^2>0$.
Minimality excludes an exceptional smooth rational $(-1)$-curve;
adjunction therefore gives $K_S\cdot E=2p_a(E)-2+b\ge0$.
The inverse of the positive matrix with nonpositive off-diagonal
entries is nonnegative by \cref{lem:stieltjes}. Consequently the
canonical pullback equations give $\lambda_i\ge0$; klt gives
$\lambda_i<1$. If $p_a(E)\ge1$, its equation
$b\lambda_E=K_S\cdot E+\sum_{E'\ne E}(E\cdot E')\lambda_{E'}$
would give $\lambda_E\ge1$. Thus every $E$ has arithmetic genus zero,
is smooth rational, and has square at most $-2$.
If the incidence multigraph had a cycle, take its vertices $J$.
Each vertex has total intersection at least two within $J$, so
$A_J\mathbf1\le q_J$. Restriction of the full discrepancy equations
and positivity of $A_J^{-1}$ would give
$\lambda_J\ge A_J^{-1}q_J\ge\mathbf1$, a contradiction.
This also excludes two intersection points or tangency between a pair.
A triple intersection would give a triangle. Thus the exceptional
divisor is an SNC forest.

Rationality follows from \cite[Lemma~5.1]{Bernasconi}: the surface $X$
is of del Pezzo type (take boundary zero), and its minimal resolution
is birational to it. The standard blowup formulas on a rational surface
give the Picard and canonical-square assertions, as well as a birational
morphism to $\PP^2$ or a Hirzebruch surface.

It remains to determine the rank of the exceptional divisor. For any ample integral $H$ on $S$,
put $a_i=H\cdot D_i$, $x=A^{-1}a$, and $M=H+\sum_i x_iD_i$.
Then $M$ is nef, $M^2=H^2+a^{\mathsf T}A^{-1}a>0$, and its exact
null locus is $D$. A divisible multiple restricts trivially to $D$
by \cref{lem:tree-picard}; \cref{thm:keel} makes it semiample.
Its Stein contraction contracts precisely the same connected curves
as $\pi$. The image of the common graph has singleton fibers over each
normal target; the projections are proper birational and quasi-finite,
hence are isomorphisms. Thus its target is $X$.
Hence $[M]\in\pi^*N^1(X)_{\QQ}$.
The map $[H]\mapsto[M]$ is the linear orthogonal projection onto
$\langle D\rangle^\perp$. Ample classes span $N^1(S)_{\QQ}$, so its
image has dimension at most $\rho(X)=1$. It contains $L\ne0$.
The independent exceptional curves therefore number $\rho(S)-1$.
\end{proof}

\begin{lemma}[Square and degree formulas]\label{lem:square-degree}
Let $B=\sum_i\lambda_iG_i$ with $A\lambda=q$, and put
\[
  H=-\bigl(K_T+B\bigr).
\]
Then
\begin{equation}\label{eq:H-square}
  H^2=K_T^2+q^{\mathsf T}A^{-1}q.
\end{equation}
If $C\not\subset G$ is an irreducible curve and $p_i=C\cdot G_i$, then
\begin{equation}\label{eq:H-degree}
  H\cdot C=-K_T\cdot C-p^{\mathsf T}\lambda.
\end{equation}
In particular, if $C$ is a $(-1)$-curve, then
\begin{equation}\label{eq:minus-one-degree}
  H\cdot C=1-p^{\mathsf T}\lambda.
\end{equation}
\end{lemma}

\begin{proof}
The relation $A\lambda=q$ gives
\[
  B^2=-\lambda^{\mathsf T}A\lambda=-\lambda^{\mathsf T}q,
  \qquad K_T\cdot B=\lambda^{\mathsf T}q.
\]
This proves \eqref{eq:H-square}.  Formula \eqref{eq:H-degree} is the definition of $H$, and adjunction gives $-K_T\cdot C=1$ for a $(-1)$-curve.
\end{proof}

\subsection{Anticanonical contractions}

\begin{theorem}[Anticanonical contraction to a rank-one klt del Pezzo surface]\label{thm:anticanonical-contraction}
Let $T$ be a smooth projective surface over an algebraically closed field of positive characteristic.  Let
\[
  G=\sum_{i=1}^rG_i
\]
be a reduced simple-normal-crossing forest of smooth rational curves.  Assume that the classes $[G_i]$ are linearly independent and that
\[
  r=\rho(T)-1.
\]
Let
\[
  B=\sum_{i=1}^r\lambda_iG_i,
  \qquad \lambda_i<1,
  \qquad H=-\bigl(K_T+B\bigr),
\]
where the $\lambda_i$ are rational.  Suppose that
\begin{enumerate}[label=\textup{(\roman*)}]
\item $H$ is nef and big;
\item an irreducible curve $C\subset T$ satisfies $H\cdot C=0$ if and only if $C$ is one of the $G_i$.
\end{enumerate}
Then a sufficiently divisible multiple of $H$ defines a projective birational morphism
\[
  f:T\longrightarrow Y
\]
with connected fibers and exceptional locus $G$.  The target $Y$ is a normal projective rank-one klt del Pezzo surface, and, after choosing compatible canonical divisors,
\begin{equation}\label{eq:crepant-contraction}
  K_T+B=f^*K_Y,
  \qquad -K_Y\ \text{is ample},
  \qquad \rho(Y)=1.
\end{equation}
\end{theorem}

\begin{proof}
Choose a positive integer $m$ such that $mH$ and $mB$ are integral.  Since $H$ is nef and big on a surface, the exceptional locus $E(mH)$ in \cref{thm:keel} is precisely the reduced union of curves $C$ with $H\cdot C=0$, hence it is $G$.  The line bundle $\cO_T(mH)$ has degree zero on every component of every connected component of $G$.  By \cref{lem:tree-picard}, its restriction to each such component is trivial.  Keel's theorem therefore shows that $mH$ is semiample.

Let $f:T\to Y$ be the Stein factorization of the morphism defined by a sufficiently divisible basepoint-free multiple of $mH$.  Since $H$ is big, $Y$ is a surface.  The equality $f_*\cO_T=\cO_Y$ implies that the induced extension of function fields is trivial, so $f$ is birational.  A curve is contracted by $f$ exactly when it has degree zero against $H$; hence $\Exc(f)=G$.  There is an ample Cartier divisor $A$ on $Y$ and a positive integer $n$ such that
\begin{equation}\label{eq:H-pullback-ample}
  nmH\sim f^*A.
\end{equation}
Indeed, after Stein factorization $A$ is the pullback of the hyperplane class from the image, and a finite pullback of an ample divisor is ample.

The pullback of $N^1(Y)_{\QQ}$ is injective and orthogonal to all
$r=\rho(T)-1$ independent exceptional curves, so $\rho(Y)\le1$.
Since $Y$ is projective, it has an ample class and $\rho(Y)\ge1$.
Therefore $\rho(Y)=1$.

It remains to identify the canonical divisor on $Y$.  Choose a nonzero rational two-form on the common function field of $T$ and $Y$, and let $K_T$ and $K_Y$ be its divisors on the two surfaces.  Then $f_*K_T=K_Y$.  From \eqref{eq:H-pullback-ample},
\[
  K_T+B=-H\sim_\QQ f^*\!\left(-\frac{1}{nm}A\right).
\]
Pushing forward gives
\[
  K_Y\sim_\QQ-\frac{1}{nm}A.
\]
Thus $K_Y$ is $\QQ$-Cartier and $-K_Y$ is ample.  Put
\[
  E=K_T+B-f^*K_Y.
\]
The divisor $E$ is exceptional and $\QQ$-linearly trivial.  The intersection matrix of the exceptional curves is negative definite by the Hodge index theorem, so an exceptional numerically trivial divisor is zero.  Hence $E=0$, proving the crepant equality in \eqref{eq:crepant-contraction}.

Finally, the simple-normal-crossing sub-pair $(T,B)$ is sub-klt
because every coefficient of $B$ is less than one. Indeed, blowing up
a point on components with coefficients $b_1,b_2$ gives exceptional
coefficient $b_1+b_2-1<1$; at a point on just one component it gives
$b_1-1<1$, and away from the boundary it gives $-1$. Repeated point
blowups therefore keep all coefficients below one. The computation is unchanged when some
coefficients are negative.  The crepant equality $K_T+B=f^*K_Y$ identifies the discrepancies of $(Y,0)$ with those of $(T,B)$ on every higher resolution.  Hence $Y$ is klt; compare the discrepancy formalism in \cite[Section~2.3]{KollarMori}.
\end{proof}

\begin{theorem}[Numerical rank-one contraction on a rational surface]\label{thm:numerical-contraction}
Let $T$ be a smooth projective rational surface over $k$, and let
\[
  G=\sum_{i=1}^rG_i
\]
be a reduced simple-normal-crossing forest of smooth rational curves whose intersection matrix is negative definite.  Assume that the classes of the $G_i$ are independent and that $r=\rho(T)-1$.  Let
\[
  B=\sum_i\beta_iG_i,
  \qquad \beta_i<1,
  \qquad H=-\bigl(K_T+B\bigr),
\]
where the coefficients are rational.  Suppose that
\[
  H\cdot G_i=0\quad\text{for every }i,
  \qquad H^2>0,
\]
and that $H\cdot C_0>0$ for one irreducible curve $C_0\not\subset G$.  Then $G$ has a projective contraction
\[
  f:T\longrightarrow Y
\]
to a normal rank-one klt del Pezzo surface, and
\[
  K_T+B=f^*K_Y.
\]
\end{theorem}

\begin{proof}
Choose an ample integral divisor $A_0$ on $T$, put
$a_i=A_0\cdot G_i$, and let $x=A_G^{-1}a$. By \cref{lem:stieltjes},
$x_i>0$. The rational divisor
\[
 M=A_0+\sum_i x_iG_i
\]
is orthogonal to $G$, has square $A_0^2+a^{\mathsf T}A_G^{-1}a>0$,
and has positive degree on every integral curve outside $G$.
Thus $M$ is nef and big with null locus exactly $G$.
Both $H$ and $M$ lie in the one-dimensional orthogonal complement
of the independent classes of $G$. Write $H\equiv cM$. The hypothesis
$H\cdot C_0>0$, together with $M\cdot C_0>0$, gives $c>0$.
Consequently $H$ itself is nef and big with null locus exactly $G$.
The conclusion, including the descent of the canonical divisor,
now follows from \cref{thm:anticanonical-contraction}.
\end{proof}

\subsection{Orthogonal projection and nonexceptional curves}

For the following identity write $f:T\to Y$ for a minimal
rank-one klt del Pezzo resolution, $G=\Exc(f)_{\rm red}$,
$A=A_G$, and $L=f^*(-K_Y)$.

\begin{lemma}[Rank-one projection identity]\label{lem:projection}
Let $P\not\subset G$ be a $(-1)$-curve, and let $p_i=P\cdot G_i$.  Then
\begin{equation}\label{eq:green}
  p^{\mathsf T}A^{-1}p
  =1+\frac{(L\cdot P)^2}{K_Y^2}>1.
\end{equation}
Moreover,
\begin{equation}\label{eq:green-charge}
  p^{\mathsf T}A^{-1}q=p^{\mathsf T}\lambda=1-L\cdot P<1.
\end{equation}
\end{lemma}

\begin{proof}
Put $u=A^{-1}p$ and
\[
  \overline P=P+\sum_i u_iG_i.
\]
Then $\overline P$ is orthogonal to every component of $G$.  Since $\rho(Y)=1$, the orthogonal complement of the exceptional span in $N^1(T)_\QQ$ is generated by $L$.  Hence
\[
  \overline P\equiv\frac{L\cdot P}{K_Y^2}L.
\]
On the other hand,
\[
  \overline P^{\,2}=-1+p^{\mathsf T}A^{-1}p.
\]
Comparing squares proves \eqref{eq:green}.  Formula \eqref{eq:green-charge} follows from \eqref{eq:minus-one-degree}.
\end{proof}

\begin{lemma}[Null curves outside the minimal exceptional divisor]\label{lem:exterior-null-curves}
Let $(T,D,L)$ be the resolution datum of a klt del Pezzo surface of
Picard number one, and let $G$ be a nef and big divisor on $T$.  If an irreducible curve $R\not\subset D$ satisfies $G\cdot R=0$, then
\[
  R^2=K_T\cdot R=-1,
  \qquad p_a(R)=0.
\]
There are only finitely many such curves, and any two of them are disjoint.
\end{lemma}

\begin{proof}
The Hodge index theorem gives $R^2<0$.  Since $R\not\subset D$,
\[
  K_T\cdot R
  =-L\cdot R-\sum_{C\subset D}\lambda_C(C\cdot R)<0.
\]
Adjunction gives
\[
  2p_a(R)-2=R^2+K_T\cdot R.
\]
The right side is a sum of two negative integers, while the left side is at least $-2$.  Hence both integers are $-1$ and $p_a(R)=0$.  The lattice $G^\perp\cap\NS(T)$ is negative definite, so it contains only finitely many integral vectors of square $-1$.  Finally, if $R_1$ and $R_2$ were two such null curves with $R_1\cdot R_2>0$, the matrix
\[
  \begin{pmatrix}-1&R_1\cdot R_2\\ R_1\cdot R_2&-1\end{pmatrix}
\]
would fail to be negative definite.  Thus $R_1\cdot R_2=0$.
\end{proof}


\section{Rational pencils and nef thresholds}
\label{sec:pencils}
A \emph{rational ruling} is a morphism $f:S\to\PP^1$ with connected
fibers and smooth rational geometric generic fiber. Its fiber class is
the divisor class of a scheme-theoretic fiber. A class in the free
Picard group of a rational surface is \emph{primitive} if it is not an
integer multiple $dG$ with $d>1$.

\subsection{Nef divisors of square zero and rational pencils}
\begin{lemma}[A nef divisor of square zero and canonical degree minus two]
\label{lem:primitive-square-zero}
Let $F$ be an integral nef divisor on a smooth projective rational
surface $T$. If $F^2=0$ and $K_T\cdot F=-2$, then $F$ is primitive
and $|F|$ is a basepoint-free complete pencil defining a generically
separable rational ruling with connected fibers.
\end{lemma}

\begin{proof}
If $F=dG$ in the torsion-free Picard lattice with $d>1$, then $d=2$
and $G^2=0$, $K_T\cdot G=-1$, contradicting the parity
$G^2-K_T\cdot G\in2\ZZ$ supplied by Riemann--Roch.
Thus $F$ is primitive. Since $(K_T-F)\cdot F=-2$,
Serre duality and Riemann--Roch give $h^0(F)\ge2$.
Write $|F|=Z+|M|$ with fixed divisor $Z$ and nonzero moving divisor
$M$. A linear system without fixed components has nef moving class.
Both $M$ and $Z$ are $F$-null. Hodge index on $F^\perp$ gives
$M^2\le0$, whereas nefness gives $M^2\ge0$; its equality case gives
$M\equiv cF$ with $c>0$. Primitivity and integrality imply
$c\in\ZZ_{>0}$. Intersection with an ample divisor and $Z\ge0$
give $c\le1$, so $c=1$ and $Z=0$.
Two suitably chosen members now have no common component and zero
intersection, so there is no base point. The image is a curve since
$F^2=0$ and the system is nonconstant.
In the Stein factorization $f:T\to B$, the normal projective base is
smooth and $H^1(B,\mathcal O_B)$ injects into $H^1(T,\mathcal O_T)=0$.
Hence $B\simeq\PP^1$. If $G_0$ is the class of a fiber, then
$F\sim dG_0$, where $d$ is the positive degree of the line bundle on
$B$ defining the complete system. Primitivity forces $d=1$;
$f_*\mathcal O_T=\mathcal O_B$ then gives $h^0(F)=2$.
The Stein condition makes $k(B)$ relatively algebraically closed in
$k(T)$. A parameter of $k(B)$ cannot be a $p$-th power in $k(T)$,
so its differential is nonzero. The field extension is separably
generated and regular. Thus the geometric generic fiber is integral.
Flatness over the smooth base and adjunction give its arithmetic genus
zero; it is therefore a smooth rational curve. This proves the ruling
assertion.
\end{proof}

\begin{lemma}[Primitive rational rulings and their fibers]
\label{lem:ruling-fibers}
Let $f:S\to\PP^1$ be a rational ruling on a smooth projective rational
surface with primitive fiber class $F$.
Every irreducible fiber is a reduced smooth rational curve.
There is a factorization of $f$ by smooth vertical $(-1)$-curve
contractions to a Hirzebruch surface. It can start with any prescribed
vertical $(-1)$-curve. Every reduced fiber of $f$ is an SNC tree of
smooth rational curves, and a vertical $(-1)$-curve has valency at most
two. If $n_t$ is the number of reduced components of the reducible fiber
at $t$, then
\begin{equation}\label{eq:primitive-ruling-relative-picard}
 \rho(S)=2+\sum_t(n_t-1).
\end{equation}
\end{lemma}

\begin{proof}
The morphism is flat, since its base is a smooth curve and $S$ is
smooth; every scheme-theoretic fiber has arithmetic genus zero.
If a fiber is $mC$ with irreducible reduced support, primitivity gives
$m=1$, and arithmetic genus zero makes $C$ smooth rational.
For a reducible connected fiber $\sum\mu_iC_i$, its intersection
with each $C_j$ gives $C_j^2<0$. The equality
$-2=K_S\cdot F=\sum\mu_i(K_S\cdot C_i)$ selects a component of
negative canonical degree. Adjunction makes it a smooth rational
$(-1)$-curve.
Contract this curve, or start with any prescribed vertical $(-1)$-curve.
The map descends because it is constant on the exceptional curve;
the descended fiber has primitive class, since its pullback is $F$.
Iterating lowers the Picard number and ends with all fibers smooth
$\PP^1$'s. The resulting smooth $\PP^1$-fibration is a projective bundle:
$\operatorname{Br}(\PP^1)=0$ over an algebraically closed field.
Splitting its rank-two bundle gives a Hirzebruch surface of Picard
number two. Reversing the contractions blows up points of smooth
fibers or their nodes; this preserves SNC fiber trees and proves the
valency assertion, including for the prescribed first contraction.
Every blowup increases both $\rho$ and $\sum_t(n_t-1)$ by one.
This proves the formula.
\end{proof}

\subsection{The nef threshold of the anticanonical pullback}
All degrees in this subsection are measured against
$L=\pi^*(-K_X)$, not against $-K_S$. The next lemma is the
nef-threshold argument specialized to the smooth resolution; its
contraction input is Tanaka's theorem.
Whenever exterior $(-1)$-curves exist, write
\[
 \ell_{\rm ext}=\min\{L\cdot P:P\text{ is an exterior }(-1)\text{-curve}\}.
\]
The minimum exists because a Cartier multiple of $L$ has positive
integral degree on every such curve. In the next lemma, $\ell=\ell_{\rm ext}$.

\begin{lemma}[The least exterior degree is the nef anticanonical threshold]
\label{lem:nef-threshold}

Let $(S,D,L)$ be the resolution datum of a klt del Pezzo surface of Picard number one
in positive characteristic with $\rho(S)>2$. Exterior $(-1)$-curves
exist, and their least positive $L$-degree $\ell$ satisfies
\[
 N:=L+\ell K_S\ \text{is nef},\qquad
 \sup\{t\ge0:L+tK_S\text{ is nef}\}=\ell.
\]
Every exterior integral curve $C$ satisfies the stronger inequality
\begin{equation}\label{eq:least-degree-canonical-bound}
 L\cdot C\ge\ell(-K_S\cdot C)\ge\ell.
\end{equation}
In particular an effective integral divisor of $L$-degree less than
$\ell$ is supported on $D$ and has $L$-degree zero. In addition,
\[
 N\cdot D_i=\ell(-2-D_i^2),\qquad
 N^2=(1-2\ell)L^2+\ell^2K_S^2\ge0.
\]

\end{lemma}

\begin{proof}
Apply \cite[Theorem~4.4]{Tanaka} over
$\operatorname{Spec}k$ to the smooth projective surface $S$ with
boundary $\Delta=0$. The base field is excellent, $S$ is
$\QQ$-factorial, and $K_S$ is Cartier, so all its hypotheses hold.
Every $K_S$-negative extremal ray contracts by a projective morphism
$f:S\to Y$ with $f_*\mathcal O_S=\mathcal O_Y$ and
$\rho(Y)=\rho(S)-1>1$. Thus $\dim Y=2$: a projective integral curve
has Picard number one, and a point has Picard number zero.
The morphism is birational and has an exceptional curve $E$.
Then $E^2<0$ and $K_S\cdot E<0$, so adjunction makes $E$ a smooth
rational $(-1)$-curve generating that ray. It is exterior because
no component of the minimal exceptional divisor has square $-1$.
Such a ray exists since $K_S\cdot L=-L^2<0$ makes $K_S$ nonnef.
Here one uses the standard compact ample slice of the closed curve
cone to detect nonnefness on an extremal ray.
A Cartier multiple of $L$ makes the positive degrees of all exterior
$(-1)$-curves discrete, so their minimum $\ell>0$ exists.
If $N=L+\ell K_S$ were nonnef, choose an $N$-negative extremal ray.
Since $L$ is nef, that ray is $K_S$-negative and hence generated by
an exterior $(-1)$-curve $E$. But
$N\cdot E=L\cdot E-\ell\ge0$, a contradiction.
For $0\le t\le\ell$, $L+tK_S$ is a convex combination of $L$ and $N$;
for $t>\ell$ it is negative on a curve attaining the minimum.
This proves the exact threshold. For an exterior integral curve,
\[
 K_S\cdot C=-L\cdot C-\sum_i\lambda_i(D_i\cdot C)<0
\]
is an integer, so nefness of $N$ proves
\eqref{eq:least-degree-canonical-bound}. The effective-divisor assertion
follows by summing components. The final identities are direct
intersections, using $L\cdot D_i=0$ and $K_S\cdot L=-L^2$.
\end{proof}


\section{Birational replacements and singular-point counts}
\label{sec:replacements}
When an exterior $(-1)$-curve meets the exceptional divisor normally in
at most two points, one of them on a unique $(-2)$-curve, removing that
$(-2)$-curve and blowing down the exterior curve is the elementary swap
of \cite[Definition~1.3]{PalkaPelkaI}. Its preservation properties are
proved in \cite[Lemma~2.9]{PalkaPelkaI}. We use the same geometric
operation in a wider numerical form, with explicit anticanonical
coefficients and exact singular-point counts.

Fix a resolution datum $(S,D,L)$ and an exterior $(-1)$-curve $P$.
Put $p_i=P\cdot D_i$ and $q_i=K_S\cdot D_i=-2-D_i^2$.
A contact at $D_i$ is called \emph{excess} if $p_i>q_i$ and
\emph{bounded} if $p_i\le q_i$. These terms refer only to the displayed
inequality, not to a transversality assumption. We do not assume that
$D+P$ is a simple-normal-crossing divisor. A replacement theorem below
will explicitly require its retained divisor to be a
simple-normal-crossing forest.

\subsection{Restrictions on intersection multiplicities}
\begin{lemma}[Elementary discrepancy bound]\label{lem:elementary-discrepancy}
Let $D_i^2=-b$ with $b\ge3$.  Then its discrepancy coefficient satisfies
\[
  \lambda_i\ge \frac{b-2}{b}.
\]
\end{lemma}

\begin{proof}
Write the positive intersection matrix with $D_i$ first:
\[
  A=\begin{pmatrix}b&-v^{\mathsf T}\\-v&M\end{pmatrix}.
\]
The Schur complement at $D_i$ is at most $b$, so
\[
  (A^{-1})_{ii}\ge \frac1b.
\]
Since the canonical-degree vector satisfies $q\ge(b-2)e_i$ and $A^{-1}$ is entrywise nonnegative,
\[
  \lambda_i=e_i^{\mathsf T}A^{-1}q
  \ge(b-2)(A^{-1})_{ii}
  \ge\frac{b-2}{b}.
\]
\end{proof}

\begin{lemma}[Existence and rigidity of an excess contact]\label{lem:excess-contact}
The contact vector $p$ has at least one excess contact.  Every nonzero bounded contact is simple.  If a higher-weight component $D_i^2=-b$, $b\ge3$, has an excess contact, then
\[
  b=3,
  \qquad P\cdot D_i=2.
\]
In particular, there is at most one higher-weight excess contact.
\end{lemma}

\begin{proof}
If $p\le q$ coordinatewise, then positivity of $A^{-1}$ and \cref{lem:projection} give
\[
  p^{\mathsf T}A^{-1}p
  \le p^{\mathsf T}A^{-1}q
  =p^{\mathsf T}\lambda
  =1-L\cdot P<1,
\]
contrary to the projection identity.

Suppose that $0<p_i\le q_i$ and $p_i\ge2$.  Then $q_i\ge p_i\ge2$, and \cref{lem:elementary-discrepancy} gives
\[
  p_i\lambda_i
  \ge \frac{p_iq_i}{q_i+2}
  \ge1,
\]
again contradicting $p^{\mathsf T}\lambda<1$.  Thus every bounded contact is simple.

Finally, if $D_i^2=-b$ is a higher-weight excess contact, then $p_i\ge b-1$.  Hence
\[
  1>p_i\lambda_i
  \ge\frac{(b-1)(b-2)}{b}.
\]
This forces $b=3$.  Since then $\lambda_i\ge 1/3$, the inequality $p_i\lambda_i<1$ excludes $p_i\ge3$; hence $p_i=2$.  Two such excess contacts would contribute at least $4/3$ to the total discrepancy sum, so at most one exists.
\end{proof}

\begin{lemma}[Discrepancy separation in one exceptional component]
\label{lem:discrepancy-separation}
Let $B_0$ and $B_n$ be distinct higher-weight vertices in one connected component of $D$.
Then
\[
  \lambda_{B_0}+\lambda_{B_n}\geq1.
\]
\end{lemma}

\begin{proof}
Let $B_0=v_0-\cdots-v_n=B_n$ be the path joining the two vertices and let $A_P$ be
its positive intersection matrix, with the original diagonal weights.  Restricting the full
discrepancy equations to the path gives
\[
  A_P\lambda_P=q_P+c,\qquad c\geq0.
\]
Since $A_P^{-1}$ is entrywise nonnegative,
\(
  \lambda_P\geq\theta:=A_P^{-1}q_P.
\)
Put $u=\mathbf1-\theta$.  For a path,
\[
  A_Pu=e_0+e_n.
\]
Let $A_0$ be the path matrix with endpoint weights $3,3$ and every interior weight $2$.
Then $A_P\succeq A_0$, and inversion reverses the Loewner order.  Since
\[
 A_0\left(\tfrac12\mathbf1\right)=e_0+e_n,
\]
we obtain
\[
 u_0+u_n=(e_0+e_n)^{\mathsf T}A_P^{-1}(e_0+e_n)\leq1.
\]
Therefore $\theta_0+\theta_n\geq1$, and the same holds for $\lambda$.
\end{proof}

\begin{lemma}[Valencies of exceptional trees]\label{lem:exceptional-valency}
Every vertex of the dual graph of $D$ has valency at most three.
\end{lemma}
\begin{proof}
At a vertex of weight $b$, let the $d$ neighbours have weights $b_j$.
They are pairwise nonadjacent. The principal star is negative definite,
so $u=\sum_j1/b_j<b$. Its discrepancy inequalities give
\[
 b\lambda\ge b-2+\sum_j\lambda_j,\qquad
 b_j\lambda_j\ge b_j-2+\lambda,
 \qquad (b-u)\lambda\ge b+d-2-2u.
\]
Since $\lambda<1$, one obtains $u>d-2$. But $u\le d/2$, so $d<4$.
\end{proof}

\subsection{Replacing one exceptional component}

Choose a component $C\subset D$ and write $D=C+D_0$.  Order the positive intersection matrix, canonical-degree vector, and contact vector as
\begin{equation}\label{eq:replacement-block-data}
  A=
  \begin{pmatrix}
    b&-v^{\mathsf T}\\
    -v&M
  \end{pmatrix},
  \qquad
  q=\binom{b-2}{q_0},
  \qquad
  p=\binom{m}{r},
\end{equation}
where $C^2=-b$ and $m=P\cdot C$.  Assume that every contact with $D_0$ is bounded.  By \cref{lem:excess-contact},
\begin{equation}\label{eq:replacement-simple-complement}
  r\in\{0,1\}^{D_0},
  \qquad q_0\ge r.
\end{equation}
Define
\begin{align}
  g&=\bigl(b-v^{\mathsf T}M^{-1}v\bigr)^{-1},
  &\theta&=M^{-1}q_0,
  &u&=M^{-1}r,\label{eq:replacement-scalars-1}\\
  \mu&=g\bigl(b-2+v^{\mathsf T}\theta\bigr),
  &h&=r^{\mathsf T}u,
  &\eta&=r^{\mathsf T}\theta,\label{eq:replacement-scalars-2}\\
  \alpha&=v^{\mathsf T}u,
  &\tau&=m+\alpha,
  &\ell&=1-\eta-\tau\mu=L\cdot P.\label{eq:replacement-scalars-3}
\end{align}
Here $\mu$ is the original discrepancy coefficient of $C$, while the original discrepancy vector on $D_0$ is
\begin{equation}\label{eq:old-restricted-discrepancy}
  \lambda_0=\theta+\mu M^{-1}v.
\end{equation}

\begin{theorem}[Replacement of one exceptional component]\label{thm:one-component-replacement}
Assume that $C$ is the only excess contact of $P$ and that
\[
  G^\dagger=P+D_0
\]
is a simple-normal-crossing forest.  Then
\[
  0\le h\le\eta<1.
\]
Put
\begin{equation}\label{eq:replacement-new-data}
  c=\frac{1-\eta}{1-h},
  \qquad
  \lambda_0^\dagger
  =\theta-c u
  =(M-rr^{\mathsf T})^{-1}(q_0-r).
\end{equation}
The divisor
\begin{equation}\label{eq:replacement-anticanonical}
  L^\dagger
  =-K_S+cP-\sum_{D_i\subset D_0}\lambda_i^\dagger D_i
\end{equation}
is nef and big and has exact null locus $G^\dagger$.  A sufficiently divisible multiple contracts $G^\dagger$ to a rank-one klt del Pezzo surface $X^\dagger$.  On the nonminimal resolution $S\to X^\dagger$ the coefficient of $P$ in the boundary convention of \cref{thm:anticanonical-contraction} is $-c$, while the coefficients on $D_0$ are $\lambda_0^\dagger$.  They satisfy
\begin{equation}\label{eq:replacement-coefficient-bounds}
  -c<1,
  \qquad
  0\le\lambda_0^\dagger\le\lambda_0<1.
\end{equation}
Moreover,
\begin{align}
  (L^\dagger)^2-L^2
  &=\frac{(1-\eta)^2}{1-h}-\frac{\mu^2}{g}>0,
  \label{eq:replacement-square-gain}\\
  L^\dagger\cdot C
  &=\tau\frac{1-\eta}{1-h}-\frac{\mu}{g}>0.
  \label{eq:replacement-deleted-degree}
\end{align}

Let $\sigma:S\to S_1$ be the blowdown of $P$.  The induced morphism $S_1\to X^\dagger$ is the minimal resolution,
with exceptional divisor the image of $D_0$. In particular
\[
  \rho(S_1)=\rho(S)-1.
\]
If $\Gamma$ is the connected component of $D$ containing $C$, put
\[
  s_C=\#\pi_0(\Gamma-C),
\]
and let $a$ be the number of connected components of $D-C$ met by the remaining contact vector $r$.  Then
\begin{equation}\label{eq:replacement-component-count}
 \#\Sing(X^\dagger)-\#\Sing(X)
 =
 \begin{cases}
   s_C-a,&r\ne0,\\
   s_C-1,&r=0.
 \end{cases}
\end{equation}
\end{theorem}

\begin{proof}
The bounded-contact inequalities and \cref{lem:excess-contact} give $q_0\ge r$, $r_i\in\{0,1\}$ and
$0\le h\le\eta\le r^{\mathsf T}\lambda_0<1$.
Thus $M-rr^{\mathsf T}$ is positive definite with nonpositive
 off-diagonal entries. The rank-one inverse formula proves
\eqref{eq:replacement-new-data}; Stieltjes positivity gives
$0\le\lambda_0^\dagger=\theta-cu\le\theta\le\lambda_0<1$.
The negative intersection matrix and canonical-degree vector of the retained forest are
\[
 A^\dagger=\begin{pmatrix}1&-r^{\mathsf T}\\-r&M\end{pmatrix},
 \qquad q^\dagger=\binom{-1}{q_0}.
\]
Its Schur complement is positive definite, and direct multiplication
sends $(-c,\lambda_0^\dagger)^{\mathsf T}$ to $q^\dagger$.
This proves the null equations and independence of all retained classes.
On the original support, block inversion gives
\begin{equation}\label{eq:replacement-old-energy}
 p^{\mathsf T}A^{-1}p=h+g\tau^2>1,\qquad
 p^{\mathsf T}\lambda=\eta+\tau\mu=1-\ell.
\end{equation}
The respective correction terms after and before the replacement are
$q_0^{\mathsf T}M^{-1}q_0+(1-\eta)^2/(1-h)$ and
$q_0^{\mathsf T}M^{-1}q_0+\mu^2/g$.
Their difference is \eqref{eq:replacement-square-gain}; after multiplication
by $g(1-h)>0$ its numerator is
\[
 g\ell^2+2g\tau\mu\ell+\mu^2(h+g\tau^2-1)>0.
\]
The deleted degree equals $\tau c-\mu/g$; its numerator after the
same multiplication is
$g\tau\ell+\mu(h+g\tau^2-1)>0$.
Here $\tau\ge m>0$, so strictness includes $\mu=0$.
Finally
\begin{equation}\label{eq:replacement-effective-difference}
 L^\dagger=L+\mu C+
       \sum_{D_i\subset D_0}(\lambda_i-\lambda_i^\dagger)D_i+cP.
\end{equation}
Every added coefficient is nonnegative. Curves outside $D\cup P$
therefore have positive degree, as does the deleted $C$; the exact
null locus is $P+D_0$. Its rank is $\rho(S)-1$, its boundary
coefficients are less than one, and its square is positive.
The SNC forest hypothesis now permits
\cref{thm:anticanonical-contraction}, proving the claimed projective klt
anticanonical contraction.

There are at most two retained contacts on $P$: each is simple and
higher-weight, has coefficient at least $1/3$, and their total original
discrepancy sum is less than one. Because $P+D_0$ is an SNC forest, these
contacts are distinct points on $P$ in different connected components
of $D_0$, and neither is an exceptional node. Blowing down $P$ therefore
preserves smoothness of the retained curves and gives at most one
transverse double point between them. The new negative intersection matrix and canonical-degree vector are
$M-rr^{\mathsf T}$ and $q_0-r$; all retained squares are at most $-2$.
The induced resolution $S_1\to X^\dagger$ is consequently minimal,
with $\rho(S_1)=\rho(S)-1$.
Deleting $C$ replaces its exceptional connected component by $s_C$ pieces.
With no retained contact, the isolated contracted $P$ gives a smooth
point, so the singular count changes by $s_C-1$.
Otherwise its blowdown joins precisely the $a$ contacted pieces,
changing that count by $s_C-a$. Every surviving component contains
only curves of square at most $-2$ and has a singular image; each therefore contributes a singular point.
\end{proof}


\subsection{Birational modifications involving an isolated $(-2)$-curve}

\begin{theorem}[The isolated-$A_1$ exchange with its two-curve contraction]
\label{thm:isolated-node-exchange}
Let $W$ be an isolated original $(-2)$-curve and let the exterior $(-1)$-curve
$P$ meet $W$ and one vertex $B$ of another exceptional connected component
$\Delta$, once each, with no other original contact.
If $B$ has weight two, $W+2P+B$ is a complete rational fiber.
Otherwise, with $A=A_\Delta$, $q=q_\Delta$, $e=e_B$, put
$h=e^{\mathsf T}A^{-1}e$ and $\eta=e^{\mathsf T}A^{-1}q$.
Then
\begin{equation}\label{eq:isolated-mixed-range}
 B^2=-3,\qquad \tfrac12<h\le\eta<1.
\end{equation}
Write $d$ for the valency of $B$ in $\Delta$; it is positive.
Blow down $P$ and then the image of $W$, obtaining $f:S\to T$.
The images of $D-W-B$ on $T$ contract to a rank-one klt del Pezzo
surface $X_0$, with minimal resolution $T$, and
\begin{equation}\label{eq:isolated-mixed-nonterminal-gain}
 K_{X_0}^2-K_X^2=2-\eta^2/h>0,\qquad
 \#\Sing(X_0)-\#\Sing(X)=d-2.
\end{equation}
Thus for $d\ge2$ the replacement does not decrease the number of
singular points and lowers the minimal-resolution Picard number by two.
If $d=1$, let $C$ be the unique neighbour of $B$.
Blow up the transverse intersection of $B_T$ and $C_T$ on $T$.
Retain the strict transform of $B_T$ as an isolated $A_1$, omit the new
exceptional $(-1)$-curve, and retain the transformed $D-W-B$.
The support contracts to a rank-one klt del Pezzo surface $X_1$ with
unchanged singular-point count, minimal-resolution Picard decrease one,
and
\begin{equation}\label{eq:isolated-mixed-terminal-gain}
 K_{X_1}^2-K_X^2
 =\frac{2\bigl(2\eta(1-\eta)+2h-1\bigr)}{4h-1}>0.
\end{equation}
\end{theorem}

\begin{proof}
The weight-two case follows from the support intersections and
\cref{lem:primitive-square-zero}. In the other case \cref{lem:projection} gives
$1/2+h>1$, and the discrepancy sum gives $\eta<1$.
If $B^2=-b$, then $\eta\ge(b-2)h$, forcing $b=3$ and $h\le\eta$.
An isolated $B$ would have $h=1/3$, so $d\ge1$.
Order $B$ first, write
\[
 A=\begin{pmatrix}3&-a^{\mathsf T}\\-a&A_0\end{pmatrix},\qquad
 q=\binom1{q_0},\qquad
 \lambda_0=A_0^{-1}q_0,\quad u=A_0^{-1}a,
 \quad s=a^{\mathsf T}\lambda_0,\quad t=a^{\mathsf T}u.
\]
Schur complementation gives
\begin{equation}\label{eq:isolated-mixed-st}
 s=\eta/h-1,\quad t=3-1/h,\quad
 q^{\mathsf T}A^{-1}q=q_0^{\mathsf T}\lambda_0+\eta^2/h,
 \quad \lambda|_{\Delta-B}=\lambda_0+\eta u.
\end{equation}
In particular $0\le s<1$, $1<t<2$, and $0\le\lambda_0<1$.
The two given points of $P$ are distinct and avoid exceptional nodes.
After the first blowdown $W$ has square $-1$ and $B$ has square $-2$;
after the second, $B_T$ is a smooth $(-1)$-curve.
The map $f$ is an isomorphism near $D-W-B$, and
\[
 K_S=f^*K_T+W+2P,\qquad f^*B_T=B+W+2P.
\]
Thus the retained matrix on $T$ is unchanged. Its anticanonical
candidate has square $L^2+2-\eta^2/h>0$ and has degree
\[
 1-s=2-\eta/h>0
\]
on the exterior curve $B_T$.
Its boundary vector is $\lambda_0$ on $\Delta-B$, with the other
exceptional components unchanged. The support is a negative-definite SNC
forest of rank $\rho(T)-1$ and all its squares are at most $-2$.
The criterion \cref{thm:numerical-contraction} gives $X_0$ and its minimal
resolution. Removing $W$ loses one connected component and deleting $B$ replaces
$\Delta$ by $d$ pieces, so the count is $d-2$.

For $d=1$ one has $a=e_C$. The point $B_T\cap C_T$ is the unchanged
original $B$--$C$ point, disjoint from the two contracted curves and
from every other retained component. Its blowup separates the two
curves: $B$ becomes an isolated $(-2)$-curve, $C$ loses one in its
square, and the new exceptional $E$ meets just these two strict
transforms, once each at distinct points.
The remaining negative intersection matrix and canonical-degree vector are $A_0+e_Ce_C^{\mathsf T}$ and
$q_0+e_C$. Their solution is
\begin{equation}\label{eq:isolated-mixed-new-discrepancy}
 \lambda_0^\dagger=\lambda_0+\varepsilon u,\qquad
 \varepsilon=\frac{1-s}{1+t}=\frac{2h-\eta}{4h-1}.
\end{equation}
Here $0<\varepsilon<\eta$, because
$\eta-\varepsilon=2h(2\eta-1)/(4h-1)>0$.
Thus every new coefficient is in $[0,1)$, by comparison with the
original $\lambda_0+\eta u$.
The new anticanonical candidate has positive degree
$1-(\lambda_0^\dagger)_C=\varepsilon$ on the exceptional curve $E$.
The square change equals
\[
 1+2s+t-\frac{(s+t)^2}{1+t}-\frac{\eta^2}{h},
\]
which is exactly \eqref{eq:isolated-mixed-terminal-gain}.
The criterion applies again. All retained squares are at most $-2$,
so this smooth surface is the minimal resolution. It has Picard
number $\rho(S)-1$, and its two local exceptional connected components are precisely the
isolated new $B$ and the connected nonempty $\Delta-B$.
No other exceptional connected component is changed.
\end{proof}


\section{Even sets of exceptional curves and the Picard lattice}
\label{sec:nodes}
An \emph{isolated exceptional node} is a connected component of $D$ consisting of
one $(-2)$-curve. If $W_1,\ldots,W_s$ are these components, put
\[
 \mathcal C_D=\ker\left(\mathbb F_2^s\longrightarrow
 \Pic(S)/2\Pic(S),\quad (\epsilon_i)\longmapsto
 \sum_i\epsilon_i[W_i]\right).
\]
This is the standard code of disjoint nodal curves
\cite[Section~2]{DolgachevMendesLopesPardini}, here restricted to the
isolated components of the exceptional divisor. It records divisibility
of their sums in $\Pic(S)$, not the existence of arbitrary classes of
square $-2$. We will also use the standard parity fact: if a sum of
$w$ disjoint $(-2)$-curves is twice a Picard class, then $4\mid w$.
This follows from Riemann--Roch parity and holds in every
characteristic; compare \cite[Section~2]{DolgachevMendesLopesPardini}.
The next argument adapts the classical branched-cover method to the
full exceptional forest of a rank-one klt del Pezzo surface. Its rank
comparison includes every unbranched exceptional component.


\begin{theorem}[Vanishing of the isolated-node code]\label{thm:no-even-nodes}
Let $X$ be a rank-one klt del Pezzo surface over an algebraically closed field of positive characteristic different from two, and let
\[
  \pi:S\longrightarrow X,
  \qquad D=\Exc(\pi)_{\mathrm{red}},
  \qquad L=\pi^*(-K_X)
\]
be its minimal resolution.  Assume that $S$ is rational, that the connected components of $D$ are trees of smooth rational curves, and that the irreducible components of $D$ are linearly independent and number $\rho(S)-1$.  Then
\[
  \mathcal C_D=0.
\]
Equivalently, no nonempty sum of isolated nodes is twice a divisor in $\Pic(S)$.
\end{theorem}

\begin{proof}
Suppose that a nonempty collection of isolated nodes satisfies
\[
  B=W_1+\cdots+W_n=2M
\]
in $\Pic(S)$.  The curves $W_i$ are pairwise disjoint.  Since the characteristic is not two, the double cover
\[
  g:Y\longrightarrow S
\]
associated with $M$ and branched along $B$ is smooth and integral: the nonempty reduced branch divisor has an odd valuation, so its defining section is not a square in the function field.  Its ramification curves $R_i$ are disjoint $(-1)$-curves, because $g^*W_i=2R_i$ and $W_i^2=-2$.  Contract them to a smooth surface $h:Y\to T$.

Adjunction gives
\[
  K_S\cdot M=0,
  \qquad M^2=-\frac n2.
\]
The standard double-cover decomposition
\[
  g_*\cO_Y=\cO_S\oplus\cO_S(-M)
\]
and Riemann--Roch on the rational surface $S$ give
\begin{equation}\label{eq:chi-double-cover}
  \chi(\cO_Y)=2-\frac n4.
\end{equation}
Thus $4\mid n$.  Since every $W_i$ is exceptional for $\pi$, one has $M\cdot L=0$, and hence
\[
  (K_S+M)\cdot L=K_S\cdot L=-K_X^2<0.
\]
The nefness of $L$ shows that $K_S+M$ is not effective.  Therefore
\[
  p_g(Y)=h^0(S,K_S)+h^0(S,K_S+M)=0.
\]
Blowing down the $R_i$ does not change $p_g$ or $q$, so \eqref{eq:chi-double-cover} yields
\begin{equation}\label{eq:qT}
  q(T)=q(Y)=\frac n4-1.
\end{equation}

The canonical divisor of $Y$ is not pseudo-effective, because
\[
  K_Y\cdot g^*L=2(K_S+M)\cdot L<0.
\]
If $K_T$ were pseudo-effective, then
\[
  K_Y=h^*K_T+\sum_iR_i
\]
would be pseudo-effective as well.  Thus $K_T$ is not pseudo-effective.  The ruledness criterion
\cite[Theorem~2.1\textup{(I)}]{CataneseLi} shows that $T$ is birationally
ruled, with the rational case included.

The double-cover and blow-down formulas give
\[
  K_Y^2=2(K_S+M)^2=2K_S^2-n,
  \qquad
  K_T^2=K_Y^2+n=2K_S^2.
\]
If a smooth surface is birational to a ruled surface over a curve of genus $q$, then
\[
  K^2+\rho=10-8q:
\]
the equality holds on a minimal ruled surface, and every blow-up decreases $K^2$ by one and increases $\rho$ by one.  It also holds for rational surfaces, starting from $\PP^2$.  Writing
\[
  r=\rk\langle D\rangle=\rho(S)-1
\]
and using $K_S^2=10-\rho(S)$ together with \eqref{eq:qT}, we obtain
\begin{equation}\label{eq:rhoT}
  \rho(T)=2(r-n).
\end{equation}

Every connected component of $D-B$ is a rational tree disjoint from the branch divisor.  The restriction of $M$ has degree zero on every irreducible component, hence is trivial by \cref{lem:tree-picard}.  Under such a trivialization, the branch section restricts to a nonzero scalar.  That scalar has a square root because the ground field is algebraically closed and the characteristic is not two.  The induced étale double cover of the tree therefore splits into two disjoint copies.  These copies are disjoint from the ramification curves and survive on $T$.  Since the components of $D-B$ are independent and span a negative-definite lattice of rank $r-n$, the two copies span a negative-definite subspace of $N^1(T)_\QQ$ of rank
\[
  2(r-n)=\rho(T),
\]
contradicting the Hodge index theorem (if $r=n$, the equality $\rho(T)=0$ already contradicts projectivity).
\end{proof}

\subsection{Divisibility forced by a lattice index}
For a nondegenerate integral lattice $\Gamma$, let
\[
 \Gamma^\vee=\{x\in\Gamma\otimes\QQ:x\cdot\Gamma\subseteq\ZZ\}
\]
be its dual lattice. Its finite discriminant group is
$\Gamma^\vee/\Gamma$, of order $|\det\Gamma|$. For a positive integer
$I$, write $v_2(I)$ for the exponent of $2$ in its prime factorization.
The Hamming weight $\operatorname{wt}(\epsilon)$ of a binary vector is
the number of its nonzero coordinates.

\begin{lemma}[The full-index obstruction from isolated nodes]
\label{lem:picard-index}
Let $S$ be a smooth projective rational surface and let
$\Gamma\subset\Pic(S)$ be a full-rank nondegenerate sublattice.
Suppose it has an orthogonal splitting
\[
 \Gamma=\Gamma_0\oplus\bigoplus_{i=1}^t\ZZ[W_i],\qquad W_i^2=-2,
\]
where the $W_i$ are distinct isolated nodes of the original exceptional
forest $D$. Put $I=[\Pic(S):\Gamma]$.
If $t>v_2(I)$, then a nonempty sum of the $W_i$ is twice a divisor
in $\Pic(S)$. In particular this inequality is impossible on the
rank-one klt del Pezzo resolution of \cref{thm:no-even-nodes}.
For $\Gamma=\langle D,P\rangle$ with $P$ an exterior $(-1)$-curve,
write $A=A_D$, $p_i=P\cdot D_i$, and $g=p^{\mathsf T}A^{-1}p$.
Then
\[
 I^2=|\det\Gamma|=\det A\,(g-1).
\]
Thus the right side must be a positive integer square; every isolated
exceptional node disjoint from $P$ may be counted among the $t$ above.
\end{lemma}

\begin{proof}
The Picard lattice of a rational surface is torsion-free and unimodular:
this holds on $\PP^2$ and every Hirzebruch surface and is preserved
by a point blowup, which adds an orthogonal $(-1)$-summand.
Hence $|\det\Gamma|=I^2$.
The integral overlattice embeds
$G=\Pic(S)/\Gamma$ into $\Gamma^\vee/\Gamma$.
Its $2$-primary subgroup has order $2^{v_2(I)}$.
The projection to $\Gamma_0^\vee/\Gamma_0$ has image of order at most
\[
 2^{v_2(|\det\Gamma_0|)}=2^{2v_2(I)-t}.
\]
If $t>v_2(I)$, this projection has a nonzero kernel element, supported
in the discriminant groups of the $W_i$ alone. It is represented modulo
$\Gamma$ by $\frac12\sum_{i\in J}W_i$ for a nonempty subset $J$.
Lifting it to $\Pic(S)$ proves the even-sum assertion.
The last determinant formula is the Schur complement of $-A$ in
\[
 Q_\Gamma=\begin{pmatrix}-A&p\\p^{\mathsf T}&-1\end{pmatrix}.
\]
The projection identity gives $g>1$, so this matrix is nondegenerate and
has the asserted determinant in absolute value.
\end{proof}
\begin{lemma}[Riemann--Roch parity and the full Picard index]
\label{lem:picard-parity}
Keep the lattice splitting of \cref{lem:picard-index} on a
smooth rational surface, and put $s=v_2(I)$. There is a binary linear
code $\mathcal K\subseteq\FF_2^t$ of dimension at least $t-s$ such that
for every word $\epsilon\in\mathcal K$,
\[
 M_\epsilon=\tfrac12\sum_i\epsilon_iW_i\in\Pic(S),\qquad
 K_S\cdot M_\epsilon=0,\qquad
 M_\epsilon^2=-\tfrac12\operatorname{wt}(\epsilon).
\]
Every word has Hamming weight divisible by four, in every characteristic.
Consequently the index data $(I,t)=(36,4)$ and $(36,3)$ are impossible
already by Riemann--Roch parity. For $(I,t)=(24,4)$ one has the explicit
forced divisibility
\[
 W_1+W_2+W_3+W_4=2M\quad\text{in }\Pic(S),\qquad
 M^2=-2,\quad K_S\cdot M=0.
\]
The last case is excluded for isolated exceptional nodes in odd positive
characteristic by \cref{thm:no-even-nodes}.
\end{lemma}

\begin{proof}
The group $G=\Pic(S)/\Gamma$ embeds in $\Gamma^\vee/\Gamma$.
Its $2$-primary subgroup $G_2$ has order $2^s$.
Let $\mathcal K$ be the kernel of its projection to
$(\Gamma_0^\vee/\Gamma_0)_2$. The latter group has order
$2^{2s-t}$, because $|\det\Gamma_0|=I^2/2^t$.
Thus $|\mathcal K|\ge2^{t-s}$. The kernel lies in the discriminant group
$\bigoplus_i\ZZ/2$ of the node summands, so it is a binary linear code.
For each word, lifting to $\Pic(S)$ and subtracting an element of
$\Gamma$ gives exactly the half-sum $M_\epsilon$ in that Picard lattice.
Adjunction gives $K_S\cdot W_i=0$, and disjointness gives the displayed
square. The standard parity property of the nodal code recalled above
makes its weight divisible by four.
For $t=3$ the only such word is zero. For $t=4$ the only possible words
are zero and $(1,1,1,1)$, so the code has dimension at most one.
For $I=36$ one has $s=2$: the lower bounds on dimension are respectively
one at $t=3$ and two at $t=4$, both impossible.
For $I=24$ one has $s=3$ and $t=4$, so the nonzero all-one word is forced.
Its square and canonical degree are those stated. For isolated exceptional nodes on a rank-one klt del Pezzo resolution in odd positive
characteristic, \cref{thm:no-even-nodes} excludes the last case.
\end{proof}


\subsection{Lattice bounds for exceptional curve configurations}
For the lattice computations below, $A_j$ denotes a chain of $j$
weight-two vertices, and $D_4$ denotes the weight-two tree with one
trivalent vertex and three leaves. The absolute determinants of
$A_1,A_2,A_3,A_4,D_4$ are $2,3,4,5,4$, respectively. For a chain,
the determinants satisfy $d_j=2d_{j-1}-d_{j-2}$, with $d_0=1,d_1=2$;
for $D_4$, eliminating its three leaves gives $2^3(2-3/2)=4$.

\begin{lemma}[An eight-curve weight-two forest has at most four connected components]
\label{lem:eight-curve-forest}
Let $(S,D,L)$ be a rational rank-one klt del Pezzo resolution in
positive characteristic different from two. Let $F\subset D$ be a
union of connected components consisting of eight weight-two curves.
Suppose an integral nondegenerate unimodular sublattice
$U\subset\Pic(S)$ is orthogonal to $F$, and $U\oplus\langle F\rangle$
has full rank in $\Pic(S)$. Then $F$ has at most four connected
components.
\end{lemma}

\begin{proof}
If $F$ has at least five components, each has at most four vertices.
Its negative definite weight-two trees are therefore $A_1,A_2,A_3,A_4$
or $D_4$, with determinants $2,3,4,5,4$, respectively.
The eight possible rank partitions with at least five parts have
root determinants
\[
\begin{array}{c|r}
F&\det F\\\hline
3A_2+2A_1&108\\
A_3+A_2+3A_1&96\\
D_4+4A_1&64\\
A_4+4A_1&80\\
2A_2+4A_1&144\\
A_3+5A_1&128\\
A_2+6A_1&192\\
8A_1&256
\end{array}
\]
For $\Gamma=U\oplus\langle F\rangle$, unimodularity of $\Pic(S)$
requires $I^2=|\det\Gamma|=\det F$.
Only $D_4+4A_1$, $2A_2+4A_1$, and $8A_1$ remain.
Their pairs $(I,t)$, where $t$ counts the isolated $A_1$ components,
are $(8,4),(12,4),(16,8)$.
In every case $t>v_2(I)$.
These are isolated nodes of $D$ and orthogonal summands of
$\Gamma$, so \cref{lem:picard-index,thm:no-even-nodes}
exclude all three rows.
\end{proof}

\begin{proposition}[The zero-adjoint three-curve configuration]
\label{prop:zero-adjoint}
Let $(S,D,L)$ be a minimal rank-one klt del Pezzo resolution in
positive characteristic different from two. Suppose $C,B_1,B_2\subset D$
are pairwise disjoint, $P\not\subset D$ is a $(-1)$-curve, and
\[
 C^2=-2,\quad B_1^2=-3,\quad B_2^2=-\beta\ (\beta\ge3),\qquad
 P\cdot C=P\cdot B_1=P\cdot B_2=1,
\]
\[
 -K_S\equiv C+B_1+B_2+2P.
\]
Then $\#\Sing(X)\le7$. 
\end{proposition}

\begin{proof}
Suppose there are at least eight singular points. Every other exceptional
component $A$ satisfies
\[
 0\le K_S\cdot A=-(C+B_1+B_2+2P)\cdot A\le0.
\]
Hence $A^2=-2$ and it is disjoint from all four displayed curves.
The three displayed exceptional components are isolated, with discrepancies
$0,1/3,(\beta-2)/\beta$. Since
$L\cdot P=(6-\beta)/(3\beta)>0$, one has $\beta=3,4,5$.
Also $K_S^2=3-\beta$, so $\rho(S)=\beta+7$.
The remaining exceptional curves form an ADE forest $F$ of rank
$\beta+3$ with at least five connected components.
The full-rank lattice $\Gamma=\langle C,B_1,B_2,P\rangle\oplus F$
has
\[
 |\det\Gamma|=(6-\beta)\det F,
\]
since the matrix of its four-curve summand is
\[
 \begin{pmatrix}-2&0&0&1\\0&-3&0&1\\0&0&-\beta&1\\1&1&1&-1\end{pmatrix},
 \qquad \det=\beta-6.
\]
Unimodularity requires $|\det\Gamma|$ to be a square.
The possible rank partitions of $F$ with at least five parts have
largest part at most four. Using the determinants $2,3,4,5,4$ of
$A_1,A_2,A_3,A_4,D_4$ leaves exactly the following rows:
\begin{center}\small
\begin{tabular}{clrrr}
\toprule
$\beta$&$F$&$I=\sqrt{|\det\Gamma|}$&$v_2(I)$&$t=\#A_1$\\\midrule
3&$A_2+4A_1$&12&2&4\\
4&$7A_1$&16&4&7\\
4&$2A_2+3A_1$&12&2&3\\
5&$8A_1$&16&4&8\\
5&$2A_2+4A_1$&12&2&4\\
5&$D_4+4A_1$&8&3&4\\\bottomrule
\end{tabular}
\end{center}
Every listed $A_1$ is an isolated exceptional node and is an orthogonal
summand of $\Gamma$. Every row has $t>v_2(I)$, contradicting
\cref{lem:picard-index,thm:no-even-nodes}.
\end{proof}


\section{Adjoint divisors and exceptional curve configurations}
\label{sec:adjoints}
An adjoint divisor here is a divisor of the form $K_S+A$, with $A$
specified in each case. The following arguments obtain effective
representatives from Riemann--Roch and Serre duality, without a
vanishing theorem. A shortest exterior curve means an exterior
$(-1)$-curve of least positive $L$-degree.

\subsection{Configurations with a nef divisor of square one}
The three square-one configurations are the following; all displayed
curves are smooth rational, $P$ is exterior and all other components
belong to $D$:
\[
\begin{array}{c|c|c}
\text{type}&A&\text{intersections}\\\hline
\mathrm{(U1)}&W+P&W^2=-2,\ P^2=-1,\ W\cdot P=2\\
\mathrm{(U2)}&U+V+P&U^2=V^2=-2,\ P^2=-1,
\ U\cdot V=U\cdot P=V\cdot P=1\\
\mathrm{(U3)}&B+2P&B^2=-3,\ P^2=-1,\ B\cdot P=2
\end{array}
\]
For (U2), require the three displayed intersections to be distinct
points. The case where they coincide is also treated by
\cref{lem:adjacent-contact-adjoint}.


\begin{lemma}[The effective adjoint remainder of a square-one configuration]
\label{lem:square-one-adjoint}
Let $(S,D,L)$ be a minimal rank-one klt del Pezzo resolution datum,
and let $A$ be a square-one divisor centered at $P\not\subset D$, with all
other displayed components of the square-one divisor in $D$.  Put $v=L^2>0$ and let
$\epsilon=0$ for \textup{(U1)}, \textup{(U2)}, and $\epsilon=1$ for
\textup{(U3)}.  There is an effective integral divisor $N_A$ such that
\[
 \begin{gathered}
 K_S+A\sim\epsilon P+N_A,\qquad L\cdot N_A=L\cdot P-v,\\
 \operatorname{Supp}(N_A)\subseteq\Null(A),\qquad
 \operatorname{Supp}(N_A)\cap\operatorname{Supp}(A)=\varnothing.
 \end{gathered}
\]
Every exterior $A$-null curve outside the displayed support
occurs in $N_A$.  If these curves are $Q_j$ and their coefficients
are $n_j$, then
\begin{equation}\label{eq:square-one-exterior-degree}
 n_j\in\ZZ_{>0},\qquad
 \sum_j n_j(L\cdot Q_j)=L\cdot P-v,\qquad
 0<L\cdot Q_j\le L\cdot P-v<L\cdot P.
\end{equation}
Moreover, $N_A=0$ if and only if there is no such exterior null curve.
\end{lemma}

\begin{proof}
The surface $S$ is rational by \cref{prop:minimal-resolution}.
Direct intersection with the three displayed supports gives
$A^2=1$, $K_S\cdot A=-1$, and respectively the supported degree
vectors $(0,1)$, $(0,0,1)$, and $(1,0)$. An exterior curve has
nonnegative intersection with the effective $A$, so $A$ is nef.
Riemann--Roch therefore gives
\[
 \chi(\mathcal O_S(K_S+A))
 =1+\tfrac12(K_S+A)\cdot A=1.
\]
Serre duality gives $h^2(K_S+A)=h^0(-A)=0$: an effective divisor
linearly equivalent to $-A$ would have negative intersection with
the nef divisor $A$.  Hence choose $Z\ge0$ with $Z\sim K_S+A$.
Since $A\cdot Z=0$, every component of $Z$ is $A$-null.
Any such component outside $\operatorname{Supp}(A)$ is disjoint
from every displayed square-one configuration component.

In \textup{(U1)}, intersecting $Z$ with the only supported null
curve $W$ gives $-2\operatorname{coeff}_W(Z)=0$.
In \textup{(U2)}, the supported coefficients solve
\[
 \begin{pmatrix}-2&1\\1&-2\end{pmatrix}
 \binom{\operatorname{coeff}_U(Z)}
       {\operatorname{coeff}_V(Z)}=\binom00,
\]
so both vanish.  In \textup{(U3)}, the only supported null curve is
$P$, and $Z\cdot P=K_S\cdot P=-1$ gives
$\operatorname{coeff}_P(Z)=1$.  Thus $Z=\epsilon P+N_A$ has the
asserted support.  Since $L\cdot K_S=-v$ and
$L\cdot A=(1+\epsilon)L\cdot P$, its $L$-degree is as stated.

For any exterior null curve $Q$ outside the square-one support,
\cref{lem:exterior-null-curves} gives $K_S\cdot Q=-1$.
It is disjoint from $P$, so $N_A\cdot Q=-1$.  Effectivity forces
$Q$ to be a component of $N_A$, with a positive integral coefficient.
All remaining components of $N_A$ lie in $D$ and have $L$-degree
zero, proving \eqref{eq:square-one-exterior-degree}.
If there is no exterior component, then $N_A$ is supported on $D$.
For every component $C$ of $N_A$,
\[
 N_A\cdot C=K_S\cdot C=-2-C^2\ge0.
\]
Consequently $N_A^2\ge0$, contrary to negative definiteness of $D$
unless $N_A=0$.  The converse follows from the occurrence assertion.
\end{proof}

\begin{proposition}[Singular-point bounds for square-one configurations]
\label{prop:square-one-bound}

Let $(S,D,L)$ be a minimal rank-one klt del Pezzo resolution datum
in positive characteristic different from two. Let $P\not\subset D$ be a smooth rational
$(-1)$-curve with $L\cdot P=\ell_{\rm ext}$.
\begin{enumerate}[label=\textup{(\roman*)}]
\item If $C\subset D$ is of weight two and $m=C\cdot P\ge2$, then
      $P$ is the only exterior null curve of $C+mP$.
\item If $P$ centers \textup{(U1)} or \textup{(U2)}, then
      $\Null(A)\setminus D=\varnothing$, $K_S+A\sim0$,
      $L\sim_{\QQ}-K_S$, $L^2=1$, and $D$ has eight weight-two
      components and at most four connected components.
\item If $A=B+2P$ has type \textup{(U3)}, then
      $\Null(A)\setminus D=\{P\}$,
      $-K_S\sim B+P$, $L\sim_{\QQ}\tfrac23B+P$, and
      $L^2=L\cdot P=1/3$. The curve $B$ is isolated in $D$;
      the other eight exceptional components have weight two and are disjoint
      from $P$, and $D$ has at most five connected components.
\end{enumerate}
In \textup{(ii)} and \textup{(iii)}, the conclusions also hold without
shortestness whenever there is no exterior $A$-null curve outside its
displayed support. 
\end{proposition}

\begin{proof}
For \textup{(i)}, $G=C+mP$ is nef and big by its displayed support.
An additional exterior null curve $Q$ is a disjoint $(-1)$-curve by
\cref{lem:exterior-null-curves}. Projecting $L$ to the negative-definite
space $G^\perp$ and then to the orthogonal pair $P,Q$ gives
\[
 (L\cdot Q)^2\le
 \frac{(L\cdot G)^2}{G^2}-L^2-(L\cdot P)^2
 =\frac{2\ell_{\rm ext}^2}{m^2-2}-L^2<\ell_{\rm ext}^2,
\]
a contradiction. 
For a square-one configuration, the effective remainder of
\cref{lem:square-one-adjoint} excludes every additional exterior
null curve by shortestness, and gives $N_A=0$.
It gives the same equality directly under the alternative hypothesis.

For \textup{(U1)} and \textup{(U2)}, $-K_S\sim A$ is nef, and for
each $D_i$ adjunction gives $0\le A\cdot D_i=2+D_i^2\le0$.
Thus all exceptional curves have square $-2$ and all discrepancy coefficients
are zero. Since $K_S^2=1$, rationality and rank one give $\rho(S)=9$
and $\#\Irr(D)=8$. The rank-one sublattice $U=\ZZ K_S$ has determinant
one and is orthogonal to $D$. Apply \cref{lem:eight-curve-forest}.

For \textup{(U3)}, $-K_S\sim B+P$.
For any other original curve $D_i$,
\[
 0\le K_S\cdot D_i=-(B+P)\cdot D_i\le0.
\]
Hence $D_i^2=-2$ and it is disjoint from both $B$ and $P$.
The discrepancy vector is $1/3$ on $B$ and zero elsewhere, giving
$L\sim_{\QQ}\frac23B+P$ and $L^2=L\cdot P=1/3$.
Here $K_S^2=(B+P)^2=0$, so $\rho(S)=10$ and $D-B$ has eight
components. The integral sublattice $U=\langle B,P\rangle$ has matrix
\[
 \begin{pmatrix}-3&2\\2&-1\end{pmatrix}
\]
and determinant $-1$. It is orthogonal to the entire $D-B$.
The rank-eight node bound gives at most four connected components there, plus the
separate connected component $B$.
\end{proof}

\subsection{Adjacent exceptional curves and multiple intersections}

\begin{lemma}[The adjacent-contact adjoint]
\label{lem:adjacent-contact-adjoint}
Work on the minimal resolution datum $(S,D,L)$.
Let $P\not\subset D$ be a smooth rational $(-1)$-curve and let
$C,H\subset D$ be distinct curves with
\[
 C^2=-b,\quad H^2=-2,\quad b\ge2,\qquad
 C\cdot H=C\cdot P=H\cdot P=1.
\]
The three intersections need not be distinct points. Then $b\le4$.
The following divisors $A_b$ are nef and satisfy
$A_b^2=-K_S\cdot A_b=1$:
\[
\begin{array}{c|c|c|c}
 b&A_b&(A_b\cdot C,A_b\cdot H,A_b\cdot P)&Z_b\\\hline
 2&C+H+P&(0,0,1)&0\\
 3&C+H+2P&(0,1,0)&P\\
 4&C+2H+3P&(1,0,0)&H+2P
\end{array}
\]
There is a unique effective integral divisor $N$ such that
\[
 K_S+A_b\sim Z_b+N,\qquad N\sim K_S+C+H+P,
 \qquad \operatorname{Supp}(N)\cap(C\cup H\cup P)=\varnothing.
\]
It is $A_b$-null and has $L$-degree $L\cdot P-L^2$.
Every exterior $A_b$-null curve outside the displayed support
occurs in $N$ with positive integral coefficient. If $N\ne0$,
its exterior support is nonempty and each exterior component $Q$
satisfies
\[
 0<L\cdot Q\le L\cdot P-L^2<L\cdot P.
\]
In particular $N\ne0$ if a further exceptional weight-two component
$T\ne C,H$ meets $P$: indeed
$N\cdot T=C\cdot T+H\cdot T+P\cdot T>0$.
If $b=2$ and $N=0$, then $\#\Sing(X)\le7$.
\end{lemma}

\begin{proof}
For $b\ge3$, the discrepancy equations give
$2\lambda_H\ge\lambda_C$ and
$b\lambda_C\ge b-2+\lambda_H$.
Their contribution to the discrepancy sum of $P$ is at least
$3(b-2)/(2b-1)$, which is at least one for $b\ge5$.
This contradicts $\sum_{D_i}\lambda_i(P\cdot D_i)=1-L\cdot P<1$.
Direct intersection gives the table and nefness, since all other
curves have nonnegative intersection with the effective $A_b$.
Riemann--Roch and Serre duality give
$\chi(K_S+A_b)=1$ and $h^2(K_S+A_b)=h^0(-A_b)=0$.
Thus choose an effective $Z\sim K_S+A_b$.
Every component of $Z$ is $A_b$-null. Outside the displayed support,
it is disjoint from that support.
For $b=2,3,4$, the supported coefficients respectively solve
\[
 \begin{pmatrix}-2&1\\1&-2\end{pmatrix}z=\binom00,
 \qquad
 \begin{pmatrix}-3&1\\1&-1\end{pmatrix}z=\binom1{-1},
 \qquad
 \begin{pmatrix}-2&1\\1&-1\end{pmatrix}z=\binom0{-1}.
\]
These give precisely $Z_b$ and the asserted remainder $N$.
Effectivity and Hodge index imply uniqueness: after removing the
common part of two representatives, effective representatives with no common component
$N_1\equiv N_2$ would give $N_1^2=N_1\cdot N_2\ge0$ in $A_b^\perp$.
An exterior null curve is a $(-1)$-curve by
\cref{lem:exterior-null-curves}; its degree against $N$ is $-1$, so it
occurs with positive integral coefficient. The $L$-degree formula
then gives the asserted bound.
If $N$ had only exceptional components, every component $B$ would satisfy
$N\cdot B=K_S\cdot B=-2-B^2\ge0$, contradicting negative definiteness
of $D$ unless $N=0$. The assertion involving $T$ follows by
intersection with the remainder class.
For $b=2$, $N=0$ gives $-K_S\sim C+H+P$ and $K_S^2=1$.
Rationality and rank one give $\#\Irr(D)=9-K_S^2=8$.
The edge $CH$ is present in the exceptional forest, so its number of
connected components is at most seven.
\end{proof}

\begin{corollary}[The only possible multiple weight-two contact at a shortest curve]
\label{cor:multiple-contact}
Keep the datum of \cref{lem:nef-threshold}, and let $P$ attain
$\ell$. If an exceptional weight-two curve $C$ has $m=P\cdot C\ge2$, then
\[
 m=2,\qquad \ell=L^2=1,\qquad L\sim_{\QQ}-K_S,
 \qquad D_i^2=-2\quad\text{for every original }D_i.
\]
In particular higher weight-two contacts do not occur. The double
contact is a type-\textup{(U1)} square-one configuration and is excluded on a
counterexample to \cref{thm:main} by \cref{prop:square-one-bound}.
\end{corollary}

\begin{proof}
The nef class $N=L+\ell K_S$ is orthogonal to both $C$ and $P$.
Their intersection matrix has determinant $2-m^2<0$, so their span
contains a positive direction. The orthogonal space of any nonzero
nef class on a smooth projective surface is negative semidefinite
by Hodge index. Consequently $N\equiv0$.
Its degrees on exceptional components give $K_S\cdot D_i=0$; hence every
$D_i$ has weight two, the discrepancy vector is zero, and $L\sim_{\QQ}-K_S$.
In particular $\ell=1$ and $v=L^2=K_S^2$ is a positive integer.
The original projection identity and the Schur bound
$(A^{-1})_{CC}\ge1/2$ now give
\[
 \frac{m^2}{2}\le p^{\mathsf T}A^{-1}p=1+\frac1v\le2.
\]
Thus $m=2$ and $v=1$. The displayed curves $C,P$ give the square-one configuration,
including a length-two intersection concentrated at one point;
\cref{prop:square-one-bound} applies.
\end{proof}


\subsection{Intersections with three disjoint exceptional $(-2)$-curves}

\begin{proposition}[An effective adjoint from three $(-2)$-contacts]
\label{prop:cubic-adjoint}
Let $C_0,C_1,C_2\subset D$ be distinct pairwise disjoint $(-2)$-curves,
and let $P$ be an exterior $(-1)$-curve with $P\cdot C_i=1$. Put
\[
 F_1=C_0+2P+C_1,\qquad F_2=C_0+2P+C_2,\qquad
 A=C_0+C_1+C_2+3P.
\]
Then $F_1,F_2$ are fiber classes of rational rulings,
$F_1\cdot F_2=2$, and $A$ is nef with $A^2=3$ and $K_S\cdot A=-3$.
There is a unique effective integral divisor $E$ such that
\[
 E\sim K_S+C_0+C_1+C_2+2P,\qquad
 \Supp(E)\cap(C_0\cup C_1\cup C_2\cup P)=\varnothing.
\]
Its support is null for both $F_1$ and $F_2$, and
\begin{equation}\label{eq:cubic-degree}
 L\cdot E=2L\cdot P-L^2.
\end{equation}
Every exterior curve $Q\ne P$ null for both rulings is a $(-1)$-curve
and occurs in $E$ with positive integral coefficient. If $E=0$, or
if there is no such $Q$, then $n(X)\le7$.
\end{proposition}
\begin{proof}
Intersections on the displayed support give $F_i^2=0$,
$K_S\cdot F_i=-2$, $F_1\cdot F_2=2$, and
\[
 A^2=3,\quad K_S\cdot A=-3,\quad
 (A\cdot C_0,A\cdot C_1,A\cdot C_2,A\cdot P)=(1,1,1,0).
\]
These divisors have nonnegative degree on every other integral curve,
so they are nef. The rational-pencil assertion follows from
\cref{lem:primitive-square-zero}. Riemann--Roch gives
$\chi(\cO_S(K_S+A))=1$, while Serre duality gives
$h^2(K_S+A)=h^0(-A)=0$. Choose $Z\ge0$ with $Z\sim K_S+A$.
Since $A\cdot Z=0$, every component of $Z$ is $A$-null.
The only such component in the displayed support is $P$; all others
are disjoint from that support. As $Z\cdot P=-1$, the coefficient
of $P$ in $Z$ is one. Thus $Z=P+E$ with $E$ as stated.

The support of $E$ is null for $F_1+F_2$, which is nef and big.
Two distinct effective representatives of its class, after subtraction
of their common components with the smaller coefficients, would give
nonzero effective divisors $E_1\equiv E_2$ with no common component
and $E_1^2=E_1\cdot E_2\ge0$, contrary to the Hodge index theorem.
This proves uniqueness. Intersecting its class with $L$ gives
\eqref{eq:cubic-degree}.

An exterior curve $Q\ne P$ null for both $F_i$ is disjoint from
$P,C_0,C_1,C_2$ and is a $(-1)$-curve by
\cref{lem:exterior-null-curves}. Hence $E\cdot Q=K_S\cdot Q=-1$,
which forces its occurrence in $E$. If there is no such exterior
curve, $E$ is supported on $D$ and every component $B$ of $E$ has
$E\cdot B=K_S\cdot B\ge0$. Negative definiteness forces $E=0$.
In that case $-K_S\sim C_0+C_1+C_2+2P$ is nef of square two.
Adjunction gives $D_i^2=-2$ for every $D_i$, so $L=-K_S$ up to
rational linear equivalence and $K_S^2=2$. Therefore
$\#\Irr(D)=9-K_S^2=7$, proving the bound.
\end{proof}

\begin{proposition}[A singular-point bound for a single exterior adjoint component]
\label{prop:three-contact-single-component}

Work on the resolution datum in positive characteristic different from two.
Let $C_0,C_1,C_2\subset D$ be disjoint weight-two curves meeting the
exterior $(-1)$-curve $P$ once each. Suppose an exterior $(-1)$-curve
$R$ is disjoint from all four curves and
\begin{equation}\label{eq:single-component-adjoint}
 K_S+C_0+C_1+C_2+2P\sim R.
\end{equation}
Then $\#\Sing(X)\le7$. 
\end{proposition}

\begin{proof}
Put $p=L\cdot P$, $v=L^2$ and $r=L\cdot R$.
The displayed identity gives $K_S^2=1$ and $r=2p-v>0$.
Suppose $\#\Sing(X)\ge8$. Rationality and rank one give exactly
$9-K_S^2=8$ exceptional components, so all eight are isolated in $D$.
Write their squares as $-b_i$, with $b_i\ge2$. Then
\[
 v=1+\sum_i\frac{(b_i-2)^2}{b_i},\qquad
 p=1-\sum_i\frac{(b_i-2)(P\cdot D_i)}{b_i}\le1.
\]
Since $v<2p\le2$, every $b_i$ is two or three, and the number $k$
of weight-three components belongs to $\{0,1,2\}$.
Let $u$ be the sum of the contacts of $P$ with those components.
Then
\[
 v=1+k/3,\qquad p=1-u/3,\qquad r=1-(k+2u)/3>0.
\]
For $k=1,2$ this forces $u=0$; for $k=0$ it is zero by definition.
Thus $p=1$ and $P$ meets no higher-weight component.
The rank-one projection identity now gives
\[
 \frac12\sum_{b_i=2}(P\cdot D_i)^2=1+\frac1v
 =\begin{cases}2&k=0,\\7/4&k=1,\\8/5&k=2.\end{cases}
\]
The left side is a half-integer, so $k=0$ and $p=v=r=1$.
The three weight-two contacts contribute three to the sum of squares, whose
value is four. Hence there is exactly one further original component $W$
with $P\cdot W=1$, and all other contacts of $P$ are zero.
Intersecting \eqref{eq:single-component-adjoint} with $W$
gives $R\cdot W=2$. Thus $W+R$ is a type-\textup{(U1)} square-one configuration.
All discrepancies are zero, so every exterior $(-1)$-curve has
$L$-degree one and $R$ is shortest.
\Cref{prop:square-one-bound} gives at most four singular
points, contrary to the assumption of eight.
\end{proof}

\begin{theorem}[Strict adjoint descent from any three weight-two contacts]
\label{thm:three-contact-descent}
Let $P,C_0,C_1,C_2$ be as in
\cref{prop:three-contact-single-component}, without assuming its adjoint identity.
Then either $\#\Sing(X)\le7$, or there is an exterior
$(-1)$-curve $Q$, disjoint from $P\cup C_0\cup C_1\cup C_2$, such that
\begin{equation}\label{eq:three-end-linear-descent}
 0<L\cdot Q\le L\cdot P-\tfrac12 L^2<L\cdot P.
\end{equation}
It may be chosen from the unique effective divisor
$E\sim K_S+C_0+C_1+C_2+2P$ of
\cref{prop:cubic-adjoint}.

\end{theorem}

\begin{proof}
Put $p=L\cdot P$ and $v=L^2$.
The proof of
\cref{prop:cubic-adjoint} gives
$E\ge0$, disjoint from the four displayed curves, of $L$-degree
$2p-v$. Its support is null for $F_1+F_2$, where
$F_i=C_0+2P+C_i$. Hodge index implies uniqueness in its linear
class, by the same common-part argument as for a square-one adjoint.
Write
\[
 E=\sum_j a_jQ_j+Z,\qquad
 a_j\in\ZZ_{>0},\quad Q_j\not\subset D,\quad Z\ge0\text{ on }D,
 \qquad M=\sum_j a_j.
\]
The $Q_j$ are $(-1)$-curves by
\cref{lem:exterior-null-curves}.
If $M=0$, then for every component $B$ of $E$,
$E\cdot B=K_S\cdot B\ge0$; negative definiteness forces $E=0$.
Then $-K_S\sim C_0+C_1+C_2+2P$ has square two, and
$\#\Irr(D)=7$, proving the first outcome.
If $M=1$, write $E=R+Z$. Since $E\cdot R=-1$ and $R^2=-1$,
$Z\cdot R=0$. Hence each component $B$ of $Z$ satisfies
$Z\cdot B=E\cdot B=K_S\cdot B\ge0$, forcing $Z=0$.
Now \cref{prop:three-contact-single-component} applies.
In all remaining cases $M\ge2$ and
\[
 \min_j L\cdot Q_j\le
 \frac{\sum_j a_j(L\cdot Q_j)}{M}
 =\frac{2p-v}{M}\le p-\frac v2.
\]
Choosing an exterior component of least degree
proves the second outcome.
\end{proof}


\section{Rational rulings with sections and bisections}
\label{sec:rulings}
Fix a rational ruling $f:S\to\PP^1$ with fiber class $F$.
A component of $D$ is \emph{vertical} if it lies in a fiber and
\emph{horizontal} otherwise. A horizontal component of degree
$F\cdot H=1$ is a section, and one of degree two is a bisection.
All fibers in the intersection equations are scheme-theoretic fibers;
their component multiplicities are retained.

A
\emph{minimal counterexample} is a klt del Pezzo surface of Picard
number one with at least eight singular points whose minimal
resolution has the smallest Picard number among all such surfaces
over the same field. Any birational replacement with at least as
many singular points and smaller minimal-resolution Picard number
contradicts this choice.

\subsection{Connected exceptional divisors and bisection ramification}
\begin{lemma}[Connected-component count from a rational ruling]
\label{lem:forest-count}
Let
\[
  f:S\longrightarrow\PP^1
\]
be a rational ruling with primitive fiber class $F$, and write
\[
  D_{\mathrm h}=H_1+\cdots+H_s,
  \qquad
  D_{\mathrm v}=D-D_{\mathrm h}
\]
for the horizontal and vertical parts of $D$.  For every reducible
fiber
\[
  \Phi_t=f^{-1}(t)_{\mathrm{red}},
\]
write
\[
  (D_{\mathrm v})_t
   :=\sum_{\substack{C\subset D_{\mathrm v}\\ f(C)=t}}C
   =R_{t,1}+\cdots+R_{t,r_t}
\]
as the decomposition into connected reduced divisors; put $r_t=0$ when
the exceptional vertical part is empty.  Define
\[
  a_t
   =\sum_{i=1}^s\sum_{j=1}^{r_t}H_i\cdot R_{t,j},
  \qquad
  q_{\mathrm h}
   =\sum_{1\le i<j\le s}H_i\cdot H_j.
\]
Then
\begin{equation}\label{eq:collapsed-fiber-euler}
  \#\pi_0(D)
   =s+\sum_t(r_t-a_t)-q_{\mathrm h}.
\end{equation}
\end{lemma}

\begin{proof}
A vertical component of $D$ lies in a reducible fiber: an irreducible
fiber has square zero and cannot belong to the negative-definite
span of $D$.  In the dual forest of $D$, contract graph-theoretically
each connected subtree carried by one $R_{t,j}$.  The resulting graph is
a forest with
\[
  s+\sum_t r_t
\]
vertices.  Its edges are exactly the horizontal--vertical incidences,
of which there are $\sum_ta_t$, and the direct
horizontal--horizontal incidences, of which there are $q_{\mathrm h}$.
The number of connected components of a forest is its number of vertices
minus its number of edges, proving
\eqref{eq:collapsed-fiber-euler}.
\end{proof}

\begin{lemma}[Ramification of a smooth bisection]
\label{lem:bisection-ramification}

Let $f:S\to\PP^1$ be a rational ruling over an algebraically closed
field of characteristic $p>0$, and let $B\simeq\PP^1$ be a smooth
bisection. Put $h=f|_B$ and
$\rho_t(B)=\sum_{x\in h^{-1}(t)}(e_x(h)-1)$, where $e_x(h)$ is the
multiplicity in the scheme-theoretic fiber of $h$.
If $p\ne2$, then $h$ is separable and tame and
\begin{equation}\label{eq:one-bisection-ramification-budget}
 \sum_t\rho_t(B)=2.
\end{equation}
For two bisections in this characteristic one therefore has
\begin{equation}\label{eq:two-bisection-ramification-budget}
 \sum_t\bigl(\rho_t(B_1)+\rho_t(B_2)\bigr)=4.
\end{equation}
If $p=2$ and $h$ is separable, it has exactly one ramification point,
with different exponent two; in this case $\sum_t\rho_t(B)=1$.
If $p=2$ and $h$ is inseparable, it is Frobenius followed by an
isomorphism, every fiber is twice one point, and the sums of these
$\rho_t(B)$ over finite sets of fibers have no uniform upper bound.
In every case, if $f^*(t)=\sum_Qm_QQ$ and $x\in B\cap f^{-1}(t)$, then
\[
 e_x(h)=\sum_{Q\ni x}m_QI_x(B,Q).
\]
Thus meeting a fiber component of multiplicity at least two forces
$e_x(h)=2$; the total number of such points is at most two when $h$
is separable, but is not so bounded in the inseparable case.
\end{lemma}

\begin{proof}
We use the different form of Riemann--Hurwitz
\cite[Tag~0C1B]{Stacks}. The finite morphism $h$ has total
degree two. Its inseparable degree is
a power of $p$. For $p\ne2$ it is therefore separable, and every local
index is one or two, prime to $p$. The residue fields are the
algebraically closed ground field, so the map is tame. Riemann--Hurwitz
makes the sum of the different exponents equal to
$-2-2(-2)=2$, and in the tame case these are $e_x-1$.
This proves both equalities.
For a separable degree-two map in characteristic two, a ramified local
parameter has the form $u=t^2a(t)$ with $a(0)\ne0$. The nonzero
differential is $t^2a'(t)\,dt$, so its order, the different exponent,
is at least two. Riemann--Hurwitz still gives total different two.
There is consequently exactly one ramification point, with that
exponent equal to two and index two.
If $h$ is inseparable, its defining rational function belongs to
$k(t^2)$ and has degree two. It is therefore $g(t^2)$ with $g$ a
fractional linear transformation. This proves the Frobenius description
and the fiber assertion; its differential is identically zero, so the
separable Riemann--Hurwitz calculation does not apply.
Finally restrict the divisor of a local parameter at $t$ from the smooth
surface to $B$. Its order at $x$ is the sum of the component
multiplicities times the local intersection numbers, proving the formula.
A summand of size at least two uses the entire degree of a fiber of $h$.
\end{proof}


\subsection{Fiber multiplicities and adjoints of bisections}

\begin{lemma}[Fibers of least possible anticanonical degree]
\label{lem:fiber-degree-equality}
Let $(S,D,L)$ be the resolution datum of a klt del Pezzo surface of Picard number one
in positive characteristic. Suppose exterior $(-1)$-curves exist and
there is a rational ruling with primitive fiber class $F$ satisfying
$L\cdot F=2\ell_{\rm ext}$. Denote by $\mu_C$ the multiplicity of a
component $C$ in its scheme-theoretic fiber. In every reducible fiber,
\[
 \sum_{R\not\subset D}\mu_R=2,\qquad
 L\cdot R=\ell_{\rm ext}\quad(R\not\subset D),
 \qquad C^2=-2\quad(C\subset D\text{ vertical}).
\]
Thus its exterior multiplicity pattern is $(2)$ or $(1,1)$.
In the first case the exterior fiber component has valency at most two and the
exceptional vertical part has at most two connected pieces.
In the second case the entire reduced fiber is a chain
\[
 R-C_1-\cdots-C_l-Q\quad(l\ge0),
\]
where the $C_i$ are retained $(-2)$-curves and every multiplicity is one.
If $s$ is the number of horizontal components of $D$, exactly $s-1$
reducible fibers have pattern $(1,1)$.
\end{lemma}

\begin{proof}
Every exterior component of a reducible fiber is a smooth rational
$(-1)$-curve, by the fiber intersection equation and adjunction with
$K_S=-L-\sum\lambda_CC$.
The least-degree inequality and $L\cdot D=0$ give
$\sum\mu_R\le2$. On the other hand adjunction against the entire fiber
gives
\[
 -2=K_S\cdot F
   =-\sum_{R\not\subset D}\mu_R
       +\sum_{C\subset D\text{ in the fiber}}\mu_C(-2-C^2).
\]
All terms in the latter sum are nonnegative. Therefore equality holds
in both estimates, every retained curve has weight two, and every omitted
curve has degree $\ell_{\rm ext}$.
The reduced fiber is an SNC tree by
\cref{lem:ruling-fibers}.
For a multiplicity-one exterior fiber component the equation $F\cdot R=0$ says
that the sum of its neighbours' multiplicities is one. It is a leaf
with a multiplicity-one neighbour. At a retained multiplicity-one
$(-2)$-curve the next neighbour is again unique and of multiplicity one.
This traces the whole connected fiber to the other exterior fiber component,
proving the chain assertion. A unique exterior fiber component of multiplicity two
has neighbours whose multiplicities sum to two, proving the other bound.
Finally, let $n_t$ be the reduced component count and $o_t$ the number
of exterior components in the fiber at $t$. Then
\[
 \rho(S)=2+\sum_t(n_t-1),\qquad
 \rho(S)-1=\#\Irr(D)=s+\sum_t(n_t-o_t).
\]
Subtracting gives $\sum_t(o_t-1)=s-1$. Each $o_t$ is one or two,
so this is the asserted exact number of two-exterior-component fibers.
\end{proof}

\begin{lemma}[A bisection adjoint is zero or one exterior curve]
\label{lem:bisection-adjoint}
Keep the ruling of \cref{lem:fiber-degree-equality}, and let
$H\subset D$ be a bisection with $H^2=-b$.
There is an effective integral divisor $N\sim K_S+F+H$, and exactly
one of the following holds:
\[
 N=0,
 \qquad\text{or}\qquad
 N=R\text{ for an exterior vertical }(-1)\text{-curve }R
       \text{ disjoint from }H.
\]
In the first case $K_S^2=4-b$, and in the second $K_S^2=3-b$.
If $N=0$, then $H$ is the only horizontal component of $D$.
If $N=R$ and another horizontal component $T$ exists, then the fiber
multiplicity of $R$ is one, and every such $T$ satisfies
\[
 T^2=-2,\quad T\cdot H=0,\quad T\cdot R=F\cdot T;
\]
its entire fiber intersection is on $R$.
\end{lemma}

\begin{proof}
Put $A=F+H$. Then $A^2=4-b$, $K_S\cdot A=b-4$, and
$\chi(K_S+A)=1$. Since $(-A)\cdot F=-2$, Serre duality gives
$h^2(K_S+A)=h^0(-A)=0$. Hence $N\ge0$ exists.
It is vertical, since $N\cdot F=0$. Its $L$-degree is
$2\ell_{\rm ext}-L^2<2\ell_{\rm ext}$.
An exterior component of an irreducible fiber already has degree
$2\ell_{\rm ext}$; all other exterior vertical components have degree
$\ell_{\rm ext}$ by \cref{lem:fiber-degree-equality}.
Thus $N$ has at most one exterior component, with coefficient one.
If it has none, it is supported on the weight-two vertical part of $D$;
for each component $C$ of $N$ one has $N\cdot C=H\cdot C\ge0$.
Negative definiteness forces $N=0$.
Otherwise write $N=R+Z$, with $Z\ge0$ on $D$.
The equality $N\cdot H=K_S\cdot H+F\cdot H+H^2=0$ shows that
$N$ is disjoint from $H$. Hence $N\cdot R=-1$, which gives
$Z\cdot R=0$. For a component $C$ of $Z$, one has
$Z\cdot C=N\cdot C=H\cdot C=0$.
Again negative definiteness gives $Z=0$.
The two identities for $K_S^2$ follow by squaring
$-K_S\sim F+H-N$, using $F\cdot N=H\cdot N=0$.
For another horizontal component $T$, adjunction gives
\[
 0\le K_S\cdot T=N\cdot T-F\cdot T-H\cdot T.
\]
This excludes $N=0$. For $N=R$, comparison with
$F\cdot T\ge\mu_R(R\cdot T)$ forces $\mu_R=1$ and equality
throughout. Thus $K_S\cdot T=H\cdot T=0$ and all of its fiber
intersection lies on $R$, as claimed.
\end{proof}

\subsection{The singular-point bound from sections and bisections}

\begin{theorem}[The singular-point bound from sections and bisections]
\label{thm:two-contact-ruling}
Let $(S,D,L)$ be the resolution datum of a minimal counterexample in positive characteristic different from two. Suppose $P$ is globally
shortest among exterior $(-1)$-curves, and
\[
 F=U+2P+V
\]
is a complete rational fiber with $U,V\subset D$ disjoint
weight-two curves meeting $P$ once. If $F\cdot C\le2$ for all
$C\in\Irr(D)$, then $\#\Sing(X)\le7$.

\end{theorem}

\begin{proof}
Put $\ell=L\cdot P$. The fiber degree equalities apply since
$L\cdot F=2\ell$. Denote the number of horizontal exceptional components by
$s$, which is positive: otherwise the independent original span of rank
$\rho(S)-1$ and the additional class $F$ would all lie in $F^\perp$.
Use the notation $r_t,a_t,q_{\rm h}$ of
\cref{lem:forest-count}.

Suppose first that all horizontal components are sections.
On the distinguished fiber each meets exactly one of $U,V$, and \cref{lem:exceptional-valency} gives $s\le6$. A one-exterior-component fiber has a multiplicity-two omitted
curve, so every section meets its retained part and $a_t=s$,
$r_t\le2$. The distinguished fiber has $r_0=2$, $a_0=s$.
A two-exterior-component fiber has $r_t\le1$, so $r_t-a_t\le1$.
For $s\ge2$, all one-exterior-component contributions other than the distinguished fiber are
nonpositive. There are $s-1$ two-exterior-component fibers, giving
\[
 \#\pi_0(D)=s+\sum_t(r_t-a_t)-q_{\rm h}
 \le s+(s-1)+(2-s)=s+1\le7.
\]
For $s=1$, every reducible fiber is a one-exterior-component fiber and the sole
section meets retained support in that fiber. These distinct fibers
use distinct edges incident with the section, so there are at most
three. Each contributes at most one, giving $\#\pi_0(D)\le4$.

Now take a bisection $H$ and its divisor $N$ from
\cref{lem:bisection-adjoint}.
If $H$ is the only horizontal component, all reducible fibers have
one exterior fiber component of multiplicity two. Outside the at most one fiber
supporting $N$, every retained curve $C$ satisfies
$H\cdot C=N\cdot C=0$.
The bisection therefore meets just the exterior fiber component, at one point
with scheme-theoretic fiber intersection two. Each such fiber gives
one ramification point of the tame degree-two map $H\to\PP^1$.
There are at most two by \cref{lem:bisection-ramification}.
If $N=R$, then on its fiber $H\cdot C=R\cdot C$ for retained $C$;
thus $H$ meets every retained connected piece of that fiber.
Those pieces belong to the connected component containing $H$.
Every other fiber contributes at most two connected components, giving
$\#\pi_0(D)\le1+2\cdot2=5$. The same bound holds when $N=0$.

Assume another horizontal component exists. Then $N=R$ and $R$ has
fiber multiplicity one. If $H^2=-2$, then $K_S^2=1$ and $D$ has
exactly eight components. Unless $D$ has no edges, the desired bound
already follows. In the edge-free case the fiber formula gives
$2=F\cdot H=2P\cdot H$, so $P$ meets the three disjoint weight-two
curves $U,V,H$ once each. This contradicts shortestness by
\cref{thm:three-contact-descent}.
Thus we may assume $H^2\le-3$.
The fiber containing $R$ is the chain
$Q-C_1-\cdots-C_l-R$ of \cref{lem:fiber-degree-equality}.
Every other horizontal component is of weight two and meets only $R$.
If $l=0$, intersecting $N=R$ with $Q$ gives $H\cdot Q=2$.
The contact-multiplicity restriction forces $H^2=-3$, so $H+2Q$ is a type
\textup{(U3)} square-one configuration at the shortest curve $Q$.
This contradicts \cref{prop:square-one-bound}.
For $l\ge1$, the identities $N=R$ give
\[
 H\cdot Q=1,\qquad H\cdot C_l=1,\qquad
 H\cdot C_i=0\ (i<l).
\]
The only original contacts of $Q$ are the weight-two curve $C_1$ and the
higher-weight $H$, once each. Moreover $C_1$ has valency exactly one
in $D$: its neighbour is $C_2$ if $l>1$, and $H$ if $l=1$.
Deleting $C_1$ makes $Q+(D-C_1)$ an SNC forest, including when the
three curves originally met at one point. The unique excess contact
of $Q$ is $C_1$.
The one-component replacement \cref{thm:one-component-replacement} has
$s_{C_1}=a=1$, hence the singular-point count is unchanged and the minimal-resolution
Picard number is strictly smaller. Its full contraction theorem therefore
gives a Picard descent, contrary to the chosen counterexample.
This exhausts the two-contact case.
\end{proof}


\section{Curve configurations on a minimal counterexample}
\label{sec:reduction}
Assume that $p>2$ and that a minimal counterexample exists in the sense
of \cref{sec:rulings}. The preceding results now give an exhaustive
reduction of the contacts of a shortest exterior $(-1)$-curve. The
conclusion concerns curves on its minimal resolution, not only a
numerically admissible graph.

\begin{theorem}[A single exterior adjoint curve for a minimal counterexample]
\label{thm:adjoint-reduction}

Let $(S,D,L)$ be the minimal resolution of a minimal counterexample in positive characteristic different from two. Choose a shortest
exterior $(-1)$-curve $P$, and put $\ell=L\cdot P$, $v=L^2$.
Then there are exceptional curves $C,B_1,B_2$ and a further original exterior
$(-1)$-curve $R$ satisfying all of the following:
\[
 C^2=-2,\qquad ( -B_1^2,-B_2^2 )=(3,\beta),\quad
 \beta\in\{3,4,5\},\qquad
 P\cdot C=P\cdot B_1=P\cdot B_2=1.
\]
These are all original contacts of $P$. The curves $C,B_1,B_2$ are
pairwise disjoint; $B_1,B_2$ lie in different exceptional connected components,
and the valency of $C$ in $D$ is zero or one. Moreover
\begin{equation}\label{eq:least-center-single-adjoint-configuration}
 R\sim K_S+C+B_1+B_2+2P,\qquad
 R\cap(C\cup B_1\cup B_2\cup P)=\varnothing,
 \qquad L\cdot R=2\ell-v\in[\ell,2\ell).
\end{equation}
In particular $0<v\le\ell\le(6-\beta)/(3\beta)$, and
\[
\begin{array}{c|c|c|c}
 \beta&K_S^2&\#\Irr(D)&\#\{\text{edges of }D\}\ \text{at most}\\\hline
 3&-1&10&2\\4&-2&11&3\\5&-3&12&4
\end{array}
\]
For every other original component $A$ one has the contact identity
\[
 R\cdot A=(-2-A^2)+(C+B_1+B_2)\cdot A.
\]
At least one edge of $D$ joins $\{C,B_1,B_2\}$ to another component.

\end{theorem}

\begin{proof}
The counterexample has $\rho(S)\ge9$.
A rational surface of this Picard number has a $(-1)$-curve, which is
exterior because $D$ is the minimal exceptional divisor.
The shortest degree exists by discreteness.
By \cref{cor:multiple-contact,prop:square-one-bound}, every weight-two contact of $P$ is simple.
The excess-contact lemma and the same square-one bound make every
higher-weight contact simple as well: any multiple one would have
weight three and contact two, giving type \textup{(U3)}.
There are at most two higher-weight contacts, since each contributes
at least $1/3$ to the discrepancy sum, which is less than one. Two in the same
original component are excluded by \cref{lem:discrepancy-separation}.
An excess contact exists, so there is a weight-two contact.
Two adjacent weight-two contacts are excluded by
\cref{lem:adjacent-contact-adjoint}: its nonzero remainder produces
a shorter curve, while its zero-remainder $b=2$ case gives the count.
Three disjoint ones are excluded by
\cref{thm:three-contact-descent}.

Suppose there are two weight-two contacts $U,V$.
Their divisor $F=U+2P+V$ is nef, square zero and of canonical
degree $-2$. Its primitivity follows from Riemann--Roch parity,
so \cref{lem:primitive-square-zero} gives its complete rational pencil.
A contacted higher-weight exceptional curve $B$ cannot meet $U$: applying
\cref{lem:adjacent-contact-adjoint} to $B,U,P$, the other end $V$
has positive degree against its remainder and forces a shorter curve.
Likewise $B$ cannot meet $V$. Hence $F\cdot B=2$.
Every other original component misses $P$, and its $F$-degree is at most
two by the SNC property of $D$. Thus every horizontal exceptional curve has degree at most two, and
\cref{thm:two-contact-ruling} excludes this case.

There is therefore exactly one weight-two contact $C$.
If $j$ is the number of higher-weight contacts, deleting $C$ leaves
$P+(D-C)$ an SNC forest: the remaining contacts are simple and lie
in distinct exceptional components. The one-component replacement changes the number of singular points by $s_C-1$ for $j=0$ and $s_C-j$ for $j>0$.
Any nonnegative change contradicts Picard minimality.
For $j=0$ the remaining possibility is one simple isolated-$A_1$
contact, whose projection quadratic form $1/2$ is impossible.
For $j=1$ the remaining possibility is the isolated-$A_1$ mixed
exchange, excluded by \cref{thm:isolated-node-exchange}.
Thus $j=2$ and $s_C\le1$; denote the other contacts by $B_1,B_2$.
The curve $C$ cannot be adjacent to $B_i$: the adjacent-contact
remainder would have positive intersection with the other $B_j$,
because $K_S\cdot B_j\ge1$ and $P\cdot B_j=1$.
That remainder is nonzero and would again give a shorter curve.
The strict discrepancy-sum inequality and $\lambda_{B_i}\ge1-2/(-B_i^2)$ give
$1/(-B_1^2)+1/(-B_2^2)>1/2$.
After relabelling, this gives exactly $(3,3),(3,4),(3,5)$.

Put $T=C+B_1+B_2+2P$.
Then $T^2=3-\beta$, $K_S\cdot T=\beta-3$, so
$\chi(K_S+T)=1$ and $h^2(K_S+T)=h^0(-T)=0$.
Choose $N\ge0$ with $N\sim K_S+T$.
It satisfies
\[
 N\cdot P=N\cdot C=N\cdot B_1=N\cdot B_2=0,
 \qquad L\cdot N=2\ell-v<2\ell.
\]
By \cref{lem:nef-threshold}, its total exterior coefficient
is zero or one. If zero, $N$ is supported on $D$ and
$N^2=K_S\cdot N\ge0$, since $N\cdot T=0$.
Negative definiteness gives $N=0$, contrary to
\cref{prop:zero-adjoint}.
Thus $N=R+Z$ for an exterior prime $R$ and $Z\ge0$ on $D$.
If $R=P$, the vanishing of $N$ against each of $C,B_1,B_2$
forces all three to occur in $Z$; then $Z\cdot P\ge3$, contrary
to $N\cdot P=-1+Z\cdot P=0$.
Hence $R\ne P$. The equality $N\cdot P=0$ now gives
$R\cdot P=Z\cdot P=0$, so none of $C,B_1,B_2$ occurs in $Z$.
Vanishing against those three curves forces both $R$ and $Z$ to be
disjoint from them. Therefore $N$ is disjoint from $\operatorname{Supp}(T)$.
Now $N\cdot R=K_S\cdot R<0$ implies $R^2<0$; adjunction makes
$R$ a $(-1)$-curve. It follows that $Z\cdot R=0$.
For every component $A$ of $Z$, one then has
$Z\cdot A=N\cdot A=K_S\cdot A\ge0$; hence $Z=0$.
This proves \eqref{eq:least-center-single-adjoint-configuration}.
Its degree lower bound is shortestness, and its upper bound follows
from the elementary discrepancies of $B_1,B_2$.
Squaring $K_S\sim R-T$ gives $K_S^2=2-\beta$.
Rationality gives $\#\Irr(D)=\beta+7$, and the forest count with
at least eight components gives at most $\beta-1$ edges.
Intersecting the adjoint identity with any other original $A$, using that
$P$ has no further original contacts, proves the contact formula.
If there were no edge from the three original core curves to the rest of
$D$, this formula would make every contact of $R$ bounded.
That contradicts \cref{lem:excess-contact}, proving the last assertion.
\end{proof}

\section{Exceptional dual graphs and the seven-point bound}
\label{sec:classification}
The reduction in \cref{thm:adjoint-reduction} bounds
the total number of edges. We first classify the required rational
intersection matrices, and then apply the result to the exceptional
curves of a minimal counterexample.

\Needspace{6\baselineskip}
\subsection{Inverse matrix entries for small weighted trees}
\begin{lemma}[Inverse matrix entries for trees with one higher-weight vertex]
\label{lem:rooted-trees}
Let a positive definite weighted tree have one vertex of weight $b\ge3$
(the root) and all other vertices of weight two. Assume that its valencies
are at most three and it has at most three edges. If $\alpha$ is the
root diagonal entry of its inverse positive matrix, the complete list is
\[
\begin{array}{c|l|c}
\text{edges}&\text{rooted tree}&\alpha\\\hline
0&[b]&1/b\\
1&[b,2]&2/(2b-1)\\
2&[b,2,2]&3/(3b-2)\\
2&[2,b,2]&1/(b-1)\\
3&[b,2,2,2]&4/(4b-3)\\
3&[2,b,2,2]&6/(6b-7)\\
3&\text{three-leaf star, root at the centre}&2/(2b-3)\\
3&\text{three-leaf star, root at a leaf}&1/(b-1).
\end{array}
\]
For the canonical-degree vector $q=(b-2)e_{\rm root}$, the root coefficient and correction
are $(b-2)\alpha$ and $(b-2)^2\alpha$.
For two disjoint weight-three rooted blocks let $\tau$ be the sum of
their root diagonal entries. With at most one or two total edges its
possible values are respectively
\[
 \mathcal S_1=\{2/3,11/15\},\qquad
 \mathcal S_2=\mathcal S_1\cup\{16/21,5/6,4/5\}.
\]
With at most three total edges, the values satisfying
$13/15<\tau\le14/15$ are exactly $29/33$ and $9/10$.
An all-weight-two component with a distinguished leaf has inverse entry
$1/2,2/3,3/4$ for zero, one, two edges, respectively; with three edges
its possible entries are $4/5$ and $1$.
\end{lemma}

\begin{proof}
With at most three edges the only unrooted trees are a vertex, a path,
and the three-leaf star. Their root positions give precisely the eight
rows. Eliminating weight-two leaves uses $a_i\leftarrow a_i-1/a_j$.
A weight-two chain of length $r$ attached to the root contributes
$r/(r+1)$ to its Schur complement. Thus the path rows and the centred
star follow immediately; for the leaf-rooted star the remaining
weight-two three-vertex path has centre inverse entry one.
The two fiber formulas follow by pairing $(b-2)e_{\rm root}$ with
the inverse. At $b=3$, the possible root inverse entries for exactly $0,1,2,3$ edges are
\[
 \{1/3\},\quad\{2/5\},\quad\{3/7,1/2\},\quad
 \{4/9,6/11,2/3,1/2\}.
\]
Adding two such entries with the stated total edge bound gives
$\mathcal S_1,\mathcal S_2$. For three edges the only new sums are
$7/9,29/33,1,29/35,9/10$, giving the indicated interval intersection.
Applying the same elimination to the weight-two leaf gives the last
assertion, including the leaf of the weight-two three-leaf star.
\end{proof}

\subsection{The ten candidate exceptional dual graphs}

\begin{lemma}[Classification of the candidate weighted forests]
\label{lem:ten-forests}
Let $\beta\in\{3,4,5\}$ and $n=\beta+7$.
Consider an $n$-vertex weighted forest with three distinguished vertices
$C,B_1,B_2$ of weights $2,3,\beta$, respectively. Assume:
\begin{enumerate}[label=\textup{(\alph*)}]
\item every other weight is two or three, and at most two other
      vertices have weight three;
\item the distinguished vertices are pairwise nonadjacent,
      $B_1,B_2$ lie in different connected components, $\deg C\le1$,
      and every vertex has valency at most three;
\item there are at most $\beta-1$ edges, and at least one edge joins
      a distinguished vertex to another vertex.
\end{enumerate}
Let $A$ have these weights on the diagonal and entries $-1$ at the
edges. Let $q=\operatorname{diag}(A)-2\mathbf1$ and let $p$ have
entries one on $C,B_1,B_2$ and zero elsewhere. Suppose $A$ is positive
 definite and, with
\[
 \lambda=A^{-1}q,\quad \ell=1-p^{\mathsf T}\lambda,\quad
 v=2-\beta+q^{\mathsf T}\lambda,\quad g=p^{\mathsf T}A^{-1}p,
\]
one has
\[
 0\le\lambda_i<1,\qquad 0<v\le\ell,\qquad v(g-1)=\ell^2.
\]
Then the forest is exactly one of the ten rows below. In the first
 table $[\cdots]$ denotes a chain, a bare numeral denotes an
undistinguished vertex of that weight, and $A_j$ denotes an
undistinguished weight-two chain of rank $j$.
\begin{center}\small
\begin{tabular}{ccll}
\toprule
Family&$\beta$&Fixed connected components&Remaining forest $Z$\\\midrule
A&3&$[C,3]+[B_1]+[B_2]$&$6A_1$ or $A_2+4A_1$\\
B&3&$[C,2,3]+[B_1]+[B_2]$&$5A_1$\\
C&3&$[C]+[B_1,2]+[B_2,2]+[3]$&$4A_1$\\
D&3&$[C,2]+[B_1]+[B_2]+2[3]$&$4A_1$ or $A_2+2A_1$\\
E&4&$[C]+[B_1,2]+[B_2]+2[3]$&$5A_1$, $A_2+3A_1$,\\
 &&&$A_3+2A_1$, or $2A_2+A_1$\\\bottomrule
\end{tabular}
\end{center}
Number multiple completions in their displayed order: A1, A2; D1, D2;
and E1 through E4. The family data are given below.  Here $\Delta=\det A(g-1)$, so the
last column is multiplied by the root determinant of the stated $Z$.
\[
\begin{array}{c|cccc}
\text{family}&\ell&v&g&\Delta/\det Z\\\hline
\mathrm A&2/15&1/15&19/15&12\\
\mathrm B&4/21&2/21&29/21&24\\
\mathrm C&1/5&2/15&13/10&45\\
\mathrm D&1/3&1/3&4/3&81\\
\mathrm E&1/10&1/15&23/20&54
\end{array}
\]
The absolute full-lattice determinants for the ten rows
are
\[
\begin{array}{c|rrrrrrrrrr}
\text{row}&\mathrm{A1}&\mathrm{A2}&\mathrm B&\mathrm C&\mathrm{D1}&\mathrm{D2}&\mathrm{E1}&\mathrm{E2}&\mathrm{E3}&\mathrm{E4}\\\hline
\Delta&768&576&768&720&1296&972&1728&1296&864&972.
\end{array}
\]
In particular no row has $\beta=5$.
\end{lemma}


\begin{proof}
Write $T_1,\ldots,T_k$ for the extra weight-three vertices, $0\le k\le2$.
Call a vertex of weight at least three higher-weight. In addition to $\ell>0$,
use the following explicit form of $v\le\ell$:
\begin{equation}\label{eq:small-forest-charge-difference}
  v-\ell=1-\beta+\lambda_C+2\lambda_{B_1}
             +(\beta-1)\lambda_{B_2}+\sum_i\lambda_{T_i}\le0.
\end{equation}
If $k=0$, then
$v=2-\beta+\lambda_{B_1}+(\beta-2)\lambda_{B_2}<0$, since
$\lambda_{B_1}+\lambda_{B_2}<1$ and $\beta-2\ge1$. Thus $k=1$ or $k=2$.

\emph{Each connected component has at most one higher-weight vertex.}
For any path, restriction of the full discrepancy equations gives its
own equations with a nonnegative additional vector. Stieltjes positivity
therefore bounds its coefficients from below by those of the path alone.
If two weight-three vertices are joined through weight-two vertices, both
endpoint coefficients are at least $1/2$: the constant vector $1/2$ is
exactly the path solution. If $B_1$ and some $T_i$ belong to one component,
take the path to the nearest other higher-weight vertex. It cannot be $B_2$.
For $\beta=3$, the two $1/2$ coefficients and
$\lambda_{B_2}\ge1/3$ make the right side of
\eqref{eq:small-forest-charge-difference} at least $1/6$.
For $\beta\ge4$ they give
$\lambda_{B_1}+\lambda_{B_2}\ge1$, contradicting $\ell>0$.
If $B_2$ is joined to a nearest $T_i$ by a path of $d\ge1$ edges, let
$D_d=2\beta d+\beta-2d+1$. On that path the coefficients at the intervening weight-two
vertices form an arithmetic progression. If the endpoint coefficients
are $u,w$, their endpoint equations are
$(\beta-1)u+(u-w)/d=\beta-2$ and
$2w+(w-u)/d=1$. Solving gives the lower bounds
\[
 \lambda_{B_2}\ge\frac{(\beta-2)(2d+1)+1}{D_d},\qquad
 \lambda_{T_i}\ge\frac{(\beta-1)(d+1)}{D_d}.
\]
Together with $\lambda_{B_1}\ge1/3$, these make $v-\ell$ at least
\[
 \frac{\beta d-\beta-d+5}{3D_d}>0.
\]
Thus neither core higher-weight vertex shares a component with another higher-weight
vertex. Finally, if the two extra higher-weight vertices share a component,
their coefficients are at least $1/2$ each; the isolated core lower
bounds make $v-\ell\ge(6-\beta)/(3\beta)>0$.

\emph{The component of $C$.}
Suppose first it contains a core higher-weight vertex. The path has $d\ge2$
edges. If that vertex is $B_1$, the endpoint sum is at least
$(d+2)/(2d+3)>1/2$; for $\beta\ge4$ the other core coefficient is at
least $1/2$, contradicting $\ell>0$.
For $B_2$ the endpoint sum is at least
$(\beta-2)(d+2)/((\beta-1)d+\beta)>2/3$ when $\beta\ge4$, again
contradicting the remaining coefficient $\lambda_{B_1}\ge1/3$.
For $\beta=3$ the edge bound forces the entire path $[C,2,B_i]$ and
leaves every other vertex isolated. Its values are
\[
 \ell=2/21,\qquad g=37/21,\qquad v=(7k-5)/21.
\]
For $k=2$ one has $v>\ell$; for $k=1$ the quantity
$v(g-1)-\ell^2=4/63$ is nonzero. This case is excluded.

Next suppose $C$ shares a component with $T_i$, at distance $d$.
The path gives $\lambda_C\ge1/(2d+3)$ and
$\lambda_{T_i}\ge(d+1)/(2d+3)$.
For $\beta=5$, $d\le4$ makes $\lambda_C\ge1/11>1/15$, whereas the
isolated core lower bounds require $\lambda_C<1/15$.
For $\beta=4,k=2$, these path bounds and the other extra coefficient
at least $1/3$ make
$v-\ell\ge(d+2)/(2d+3)-1/2>0$.
For $\beta=4,k=1$, the fact that $C$ is a leaf gives
$\lambda_{T_i}\le(d+1)\lambda_C\le4\lambda_C$:
along the weight-two path its coefficient increments are nonincreasing,
since every side branch contributes a nonnegative term.
Then $\ell>0$ implies
\[
 v=-2+\lambda_{B_1}+2\lambda_{B_2}+\lambda_{T_i}
 <2-3\lambda_{B_1}-2\lambda_{B_2}\le0.
\]
For $\beta=3,k=2$, the same path lower bounds give
$v-\ell\ge(d+2)/(2d+3)-1/3>0$.
It remains $\beta=3,k=1$ with at most two edges.
The only possible component of $C$ is
$[C,3]$, $[C,2,3]$, or $[C,3,2]$.
The third has $(\ell,v)=(1/12,1/6)$ and is excluded.
For $[C,3]$ the core vertices are isolated: attaching the possible second
edge to either $B_i$ would increase $v$ and decrease $\ell$ by $1/15$,
violating $v\le\ell$.
The remaining edge is therefore an unrelated weight-two edge, if present.
This gives precisely family A. The component $[C,2,3]$ uses both edges
and gives family B. Direct inversion gives their displayed numerical data.

For all remaining cases the three distinguished vertices lie in distinct
components, and $C$'s component has weight two. Put
\[
 x=(A^{-1})_{B_1B_1}-1/3,\quad
 y=(A^{-1})_{B_2B_2}-1/\beta,\quad z=(A^{-1})_{CC},\quad
 \tau=\sum_i(A^{-1})_{T_iT_i}.
\]
These quantities come from the separate inverse-matrix blocks.
In particular $x,y\ge0$, $z\ge1/2$, and $\tau\ge k/3$.
The three scalar identities become
\begin{align}
 \ell&=2/\beta-1/3-x-(\beta-2)y,\label{eq:small-forest-block-length}\\
 v&=-5/3+4/\beta+x+(\beta-2)^2y+\tau,\label{eq:small-forest-block-square}\\
 g&=z+1/3+1/\beta+x+y.\label{eq:small-forest-block-energy}
\end{align}
The inequality $v\le\ell$ is
\begin{equation}\label{eq:small-forest-block-cap}
 \tau+2x+(\beta-2)(\beta-1)y\le4/3-2/\beta.
\end{equation}

\emph{$\beta=3$.}
For $k=2$, \eqref{eq:small-forest-block-cap} forces
$x=y=0$, $\tau=2/3$. All four higher-weight vertices are isolated.
Then $v=\ell=1/3$, and the projection identity forces $z=2/3$.
Thus $C$'s component is exactly $[C,2]$; the possible remaining edge is
an unrelated weight-two edge. These are the two rows of family D.
For $k=1$, put $s=x+y$ and $u=\tau-1/3$. The projection identity becomes
\begin{equation}\label{eq:beta-three-scalar-identity}
 s(z+1/3)+u(z-1/3+s)=1/9.
\end{equation}
If $C$ is isolated, $z=1/2$. For $u=0$, this forces $s=2/15$.
With at most two edges at the components containing $B_1$ and $B_2$ the possible
positive sums $s$ are $1/15,2/21,1/6,2/15$; only one weight-two leaf
at each $B_i$ gives the required value. This is family C.
For $u>0$, the component containing the extra weight-three vertex has an edge; the core-boundary
condition requires another edge at a $B_i$. Hence $s=u=1/15$, and the
left side of \eqref{eq:beta-three-scalar-identity} is $16/225\ne1/9$.
If $C$ is not isolated, positivity of $v=s+u$ requires an edge at a component containing a higher-weight vertex as well. The two-edge bound therefore forces
$z=2/3$ and $(s,u)=(1/15,0)$ or $(0,1/15)$.
The left side is then $1/15$ or $1/45$, neither of which is $1/9$.
This finishes $\beta=3$.

\emph{$\beta=4$.}
For $k=1$, the discrepancy-sum inequality $x+2y<1/6$ makes each distinguished higher-weight
component a path with its higher-weight vertex at an end. Indeed every other
root placement in the at-most-three-edge table has root inverse entry at least
$1/2$ for $B_1$, or at least $1/3$ for $B_2$, contradicting that inequality.
Let their numbers of edges be $i,j$, so
\[
 x=\frac{i}{3(2i+3)},\qquad y=\frac{j}{4(3j+4)}.
\]
For $d=i+j\le3$, maximizing over its $d+1$ possibilities subject to
$x+2y<1/6$ gives the first row below. The second bounds the remaining
inverse entry at the extra weight-three vertex using at most $3-d$ edges:
\[
\begin{array}{c|cccc}
 d&0&1&2&3\\\hline
 \max(x+4y)&0&1/7&22/105&3/13\\
 \max\tau&2/3&1/2&2/5&1/3
\end{array}
\]
Each column sum is at most $2/3$, so
$v=-2/3+x+4y+\tau\le0$. This excludes $k=1$.
For $k=2$, \eqref{eq:small-forest-block-cap} gives $2x+6y\le1/6$.
An edge at $B_2$ gives $y\ge1/28$, which is impossible.
Two edges in $B_1$'s component give $x\ge2/21$, also impossible.
Thus $B_2$ is isolated and $B_1$ is isolated or has one weight-two leaf.
In the latter case $x=1/15$, so $\tau\le7/10$.
An edge at a component containing an extra weight-three vertex would give $\tau\ge11/15$;
hence both extras are isolated. Now $(\ell,v)=(1/10,1/15)$ and the
projection identity forces $z=1/2$. The five remaining weight-two vertices
have at most two edges, giving exactly the four completions of family E.
If both core higher-weight vertices are isolated, the required core boundary
edge is in the weight-two component of $C$.
With one edge there, $z=2/3$ and the projection identity forces $\tau=7/9$, which is not
in $\mathcal S_2$. With two edges, $z=3/4$ and the projection identity forces
$\tau=3/4\notin\mathcal S_1$.
With three edges no edge remains at a component containing an extra weight-three vertex, and
$v=0$. There are no further rows.

\emph{$\beta=5$.}
The positive length bound is $x+3y<1/15$.
Even one edge at $B_1$ gives $x\ge1/15$; one edge at $B_2$ gives
$3y\ge1/15$. Thus both are isolated and the component of $C$ has an edge.
For $k=1$, the component containing the extra weight-three vertex has at most three edges and
$\tau\le2/3$, so $v=\tau-13/15<0$.
For $k=2$, if the component of $C$ has at least two edges, the two components containing $T_1,T_2$ have at most two total edges, giving
$\tau\le5/6<13/15$ and again $v<0$.
The sole remaining case has $[C,2]$, hence $z=2/3$, $\ell=1/15$ and
$g=6/5$. The projection identity requires $v=1/45$ and $\tau=8/9$.
But $0<v\le\ell$ requires $13/15<\tau\le14/15$, and the only such
values with at most three remaining edges are $29/33$ and $9/10$,
neither of which is $8/9$. Thus $\beta=5$ is impossible.

This determines the fixed components of all five families.
For family A there are six remaining weight-two vertices and at most
one remaining edge, giving $6A_1$ or $A_2+4A_1$. Family B leaves five
isolated vertices, and family C leaves four isolated vertices.
Family D leaves four weight-two vertices with at most one edge,
giving $4A_1$ or $A_2+2A_1$. Family E leaves five weight-two vertices
with at most two edges; their edge components are empty, a single
edge, a three-vertex path, or two disjoint edges. This gives precisely
the four displayed completions. No other connected shape can occur
with these edge allowances. Inverting their fixed blocks gives the stated $\ell,v,g$ and
$\Delta/\det Z$; weight-two components of $Z$ contribute only their
usual root determinants. This proves the full ten-row classification.
\end{proof}


\subsection{Determinant, parity, and double-cover obstructions}

\begin{theorem}[The single exterior adjoint configuration is impossible]
\label{thm:forest-exclusion}
The configuration of
\cref{thm:adjoint-reduction} cannot occur on a
minimal counterexample in positive characteristic different from two.
All three values $\beta=3,4,5$ are excluded.
\end{theorem}

\begin{proof}
Retain all the exceptional curves and put $\ell=L\cdot P$ and $v=L^2$.
For another original component $A$ of weight $b=-A^2$, let
$u=(C+B_1+B_2)\cdot A\ge0$. The adjoint identity gives
\[
 R\cdot A=b-2+u.
\]
The elementary discrepancy bound and the strict discrepancy-sum inequality of the actual
$(-1)$-curve $R$ therefore imply
\[
 1>\sum_{D_i}\lambda_i(R\cdot D_i)
       \ge\frac{(b-2)^2}{b}.
\]
For $b\ge4$ the last expression is at least one, a contradiction.
Thus every additional original weight is two or three. Every additional
weight-three vertex contributes at least $1/3$ to that same sum,
so there are at most two of them. The valency-three bound is \cref{lem:exceptional-valency}. All other graph requirements of
\cref{lem:ten-forests} are conclusions of
\cref{thm:adjoint-reduction}.

For the full negative intersection matrix $A_D$, its canonical-degree
vector $q$, and the contact vector $p$ of $P$, the original identities
are
\[
 \lambda=A_D^{-1}q,\quad
 \ell=1-p^{\mathsf T}\lambda,\quad
 v=K_S^2+q^{\mathsf T}\lambda=2-\beta+q^{\mathsf T}\lambda,
 \quad v(p^{\mathsf T}A_D^{-1}p-1)=\ell^2.
\]
Here $0<v\le\ell$ is already proved by shortestness of $P$ and the
degree $L\cdot R=2\ell-v$.
Thus the complete finite matrix table applies, without assuming the
additional curves have weight two or assuming the table realizable.
The sublattice $\Gamma=\langle D,P\rangle\subset\Pic(S)$ has full rank
and determinant $\Delta$ as in that table. Unimodularity rejects every
row whose $\Delta$ is not a square. Only A2, D1, and E2 remain:
\[
\begin{array}{c|rrrr}
\text{row}&\Delta&I=\sqrt\Delta&v_2(I)&t\\\hline
\mathrm{A2}&576&24&3&4\\
\mathrm{D1}&1296&36&2&4\\
\mathrm{E2}&1296&36&2&3
\end{array}
\]
The $t$ indicated vertices are isolated weight-two components of
$D$ disjoint from $P$, hence orthogonal node summands of $\Gamma$.
By \cref{lem:picard-parity}, D1 and E2 violate Riemann--Roch parity
on any smooth rational surface. In A2, writing its four isolated nodes
as $W_1,\ldots,W_4$, the explicit forced class is
$M=(W_1+W_2+W_3+W_4)/2\in\Pic(S)$ with $M^2=-2$ and $K_S\cdot M=0$.
It contradicts \cref{thm:no-even-nodes} in odd positive characteristic.
This excludes every row.
\end{proof}

\subsection{Proof of the seven-point bound}
\begin{proof}[Proof of \cref{thm:main}]
Fix an algebraically closed field $k$ of characteristic $p>2$ and
suppose that there is a counterexample. Choose one whose minimal
resolution $S$ has the least Picard number. By
\cref{prop:minimal-resolution}, $S$ is rational and its exceptional divisor
$D$ is a simple-normal-crossing rational forest with
$\#\Irr(D)=\rho(S)-1\ge8$. In particular $\rho(S)>2$.
The nef-threshold lemma, \cref{lem:nef-threshold}, gives a
shortest exterior $(-1)$-curve and controls the degree of every
exterior integral curve.

All replacements used in \cref{thm:adjoint-reduction}
are defined over $k$. Their projectivity and klt anticanonical
contractions follow from \cref{thm:anticanonical-contraction,thm:numerical-contraction};
whenever they do not decrease the number of singular points, they
contradict the minimality of $\rho(S)$. The adjoint arguments are
intersection and Riemann--Roch arguments on this rational surface.
The discrepancy coefficients and matrix inversions are computed over
$\QQ$; their denominators impose no further condition on the
characteristic of $k$.
The bisection argument and the exceptional-node double cover are
valid because $p\ne2$. Therefore
\cref{thm:adjoint-reduction} gives the stated single
exterior adjoint configuration. Its exclusion in
\cref{thm:forest-exclusion} is a contradiction.
Thus $n(X)\le7$.
\end{proof}


\section{Frobenius constructions and sharpness}
\label{sec:examples}

The characteristic-two family is due to Keel--McKernan
\cite[Section~9]{KeelMcKernan}, and the characteristic-three seven-point
surface is Bernasconi's \cite[Section~3.1 and Proposition~3.2]{Bernasconi};
see also \cite[Example~7.4]{Lacini}. The Frobenius graph and the
successive contact blowups below recover Bernasconi's construction
exactly at $(p,n)=(3,3)$. Bernasconi also observes the higher-characteristic change of canonical
sign \cite[Remark~3.3]{Bernasconi}.

Bernasconi contracts the configuration by a surface minimal model
program; here an explicit nef divisor and its exact null locus give
the contraction directly by Keel's theorem. The calculation treats
arbitrary $p,n$, and supplies the Picard classes needed for the
subsequent analysis of degree-two maps and shortest-curve replacements.
It does not use the upper-bound theorem.

\begin{proposition}[Canonical divisors in the Frobenius family]
\label{prop:frobenius-family}

Let $k$ be algebraically closed of characteristic $p>0$, and let
$B_0\subset\PP^1_x\times\PP^1_y$ be the graph $y=x^p$.
Write $a,b$ for the fiber classes $x=\text{constant}$ and
$y=\text{constant}$. Choose $n\ge3$ distinct points of the graph.
At each one blow up $p$ times along its contact with the corresponding
$b$-fiber, obtaining $S_{p,n}$ with total-transform classes $E_{ij}$.
The strict transforms and exceptional curves have the following classes:
\[
 B=pa+b-\sum_{i,j}E_{ij},\quad
 F_i=b-\sum_{j=1}^pE_{ij},\quad
 C_{ij}=E_{ij}-E_{i,j+1}\ (1\le j<p),\quad P_i=E_{ip}.
\]
The reduced divisor $D=B+\sum_i(F_i+\sum_{j<p}C_{ij})$ consists of
one isolated curve of square $-p(n-2)$, $n$ isolated curves of square
$-p$, and $n$ disjoint weight-two chains $A_{p-1}$.
It contracts projectively to a normal klt surface $X_{p,n}$ with
$\rho(X_{p,n})=1$ and $\#\Sing(X_{p,n})=2n+1$. Denote this
contraction by $\pi:S_{p,n}\to X_{p,n}$. Put
\[
 M=B+(n-2)b,
 \qquad d_{p,n}=2-(p-2)(n-2),
 \qquad t_{p,n}=\frac{d_{p,n}}{p(n-2)}.
\]
Then $M$ is nef and big with exact null locus $D$, and
\[
 M^2=p(n-2),\qquad
 L:=\pi^*(-K_{X_{p,n}})\sim_{\QQ}t_{p,n}M,
 \qquad L^2=\frac{d_{p,n}^{\,2}}{p(n-2)},\quad L\cdot P_i=t_{p,n}.
\]
The coefficients in the canonical pullback on $B,F_i,C_{ij}$ are, respectively,
\[
 1-\frac{2}{p(n-2)},\qquad 1-\frac2p,\qquad 0.
\]
In particular $-K_{X_{p,n}}$ is ample exactly for $p=2$ (all $n\ge3$)
or for $(p,n)=(3,3)$; $K_{X_{3,4}}\sim_{\QQ}0$; and $K_{X_{p,n}}$
is ample in all other cases with prime $p$ and $n\ge3$.
For $p=3$ we also write $S_n,X_n,D_n,M_n,t_n,L_n$ and set
$U_i=C_{i1}$, $V_i=C_{i2}$, recovering the seven-point notation below.
\end{proposition}

\begin{proof}
At a finite selected point $(c,c^p)$ put $u=x-c$, $v=y-c^p$.
The graph is $v=u^p$. In the successive charts $v=u^jv_j$ its strict
transform is $v_j=u^{p-j}$ and the strict fiber is $v_j=0$.
For $j<p$ they meet at the next blowup center; at $j=p$ the graph,
fiber, and previous exceptional curve meet $P_i$ at $v_p=1,0,\infty$,
respectively. Earlier exceptional curves form the displayed chain.
Reciprocal coordinates $u=1/x$, $v=1/y$ give the same equations at
infinity. Thus the strict curves are smooth and the only local incidences
are
\[
 B\cdot P_i=F_i\cdot P_i=C_{i,p-1}\cdot P_i=1,\qquad
 C_{ij}\cdot C_{i,j+1}=1.
\]
In particular all blocks of $D$ are disjoint negative definite rational
trees. There are $np+1=\rho(S_{p,n})-1$ independent exceptional components.
The total-transform basis gives
$K_{S_{p,n}}=-2a-2b+\sum E_{ij}$.
Let $\lambda_B=1-2/[p(n-2)]$ and $\lambda_F=1-2/p$.
Solving the canonical intersection equations on these disjoint blocks gives
these coefficients on $B,F_i$, and zero on all $C_{ij}$. All belong to
$[0,1)$. Define the rational divisor
\[
 \Lambda=-\left(K_{S_{p,n}}+\lambda_BB+\lambda_F\sum_iF_i\right).
\]
Substitution of the displayed classes gives
$\Lambda=t_{p,n}M$ in $\Pic(S_{p,n})_{\QQ}$.

Direct intersection gives $M\cdot B=0$ and $M^2=p(n-2)>0$.
If an irreducible curve $R\ne B$ is not vertical for the $b$-ruling,
then $M\cdot R=B\cdot R+(n-2)b\cdot R>0$.
A nonspecial vertical fiber has degree $p$.
A special scheme-theoretic fiber is
\[
 F_i+\sum_{j=1}^{p-1}jC_{ij}+pP_i.
\]
Here $M\cdot P_i=1$ and its other components have degree zero.
This proves nefness and the exact null locus.
The restriction to each original rational tree is trivial, so
\cref{thm:keel,lem:tree-picard} give a projective Stein contraction
$\pi:S_{p,n}\to X_{p,n}$ with this exact exceptional locus.
It is birational since $M$ is big and the Stein map has equal function
fields. Its target is normal; the independent exceptional rank and an
ample target class give Picard number one. Write
$M\sim_{\QQ}\pi^*A$ with $A$ ample.
Pushforward of the canonical identity gives
$K_{X_{p,n}}\sim_{\QQ}-t_{p,n}A$, hence $K_{X_{p,n}}$ is $\QQ$-Cartier.
The divisor
$K_{S_{p,n}}+\lambda_BB+\lambda_F\sum_iF_i-\pi^*K_{X_{p,n}}$
is exceptional and numerically trivial, hence zero by negative
definiteness. Therefore $L=\pi^*(-K_{X_{p,n}})=\Lambda$, proving the
stated formulas.
The SNC pullback with coefficients less than one proves that the target
is klt, in every characteristic, without a tame-quotient assertion.
No retained curve has square $-1$, so this is the minimal resolution.
Each of its $2n+1$ connected blocks contracts to a singular point: the
exceptional locus of a nontrivial resolution over a smooth point would
contain a $(-1)$-curve. Distinct blocks have distinct images by connected
fibers. The sign of $d_{p,n}$ gives all three canonical-sign cases, since
$(p-2)(n-2)<2$ holds exactly in the positive cases stated.
\end{proof}

\subsection{Even sets and inseparable bisections in characteristic two}
The code that occurs below is the classical doubled-even code
\cite[Section~2]{DolgachevMendesLopesPardini}
\[
 DE(n)=\{(x_1,x_1,\ldots,x_n,x_n):
             (x_1,\ldots,x_n)\in\FF_2^n,\ \sum_i x_i=0\}.
\]
It has dimension $n-1$ and minimum nonzero weight four for $n\ge2$.
The code of the exceptional curves $F_i,U_i$ below is $DE(n)$.
The same surface has an inseparable bisection.

\begin{theorem}[Rulings and divisibility on the Keel--McKernan family]
\label{thm:characteristic-two}
For every algebraically closed field of characteristic two and every
$n\ge3$, the preceding construction gives a rank-one klt del Pezzo
surface $X_{2,n}$ with $2n+1$ distinct singular points and
\[
 K_{X_{2,n}}^2=\frac2{n-2}.
\]
Its minimal exceptional forest is $[2n-4]+2nA_1$.
Write $U_i=E_{i1}-E_{i2}$ and $P_i=E_{i2}$. On its original resolution,
\[
 L=\frac{2a+(n-1)b-\sum_i(E_{i1}+E_{i2})}{n-2},\quad
 \ell_{\rm ext}=\frac1{n-2},\quad
 L+\ell_{\rm ext}K_S=\frac{n-3}{n-2}b.
\]
The rational ruling of class $b$ has exactly the $n$ reducible
fibers
\[
 F_i+U_i+2P_i,
\]
whose exterior fiber components $P_i$ are shortest. Its sole exceptional
horizontal component is the purely inseparable bisection $B$.
For $n\ge4$, the complete code of two-divisible subsets of the $2n$
isolated exceptional nodes is
\[
 \mathcal C_D=
 \left\{\sum_{i\in J}(F_i+U_i):J\subseteq\{1,\ldots,n\},\ |J|
       \text{ even}\right\},
 \qquad \dim_{\FF_2}\mathcal C_D=n-1.
\]
Here the displayed effective divisor has half-class
$\frac{|J|}{2}b-\sum_{i\in J}P_i$. Its nonzero minimum support size is
four. For $n=3$, all seven exceptional curves are isolated nodes and the code
has dimension three, with all seven nonzero words of weight four.
In particular the seven-point bound is false in characteristic two,
and the isolated-node vanishing theorem does not extend to that
characteristic.
\end{theorem}

\begin{proof}
Set $p=2$ in \cref{prop:frobenius-family}; this proves the
construction, contraction, klt condition, Picard number, ample
anticanonical divisor, singular-point count, and square.
The class $(n-2)L$ is integral and has positive integral degree on every
exterior prime. Each $P_i$ has degree $1/(n-2)$, proving the
minimum without the nef-threshold lemma or the upper-bound theorem.
Substituting $K_S=-2a-2b+\sum_i(E_{i1}+E_{i2})$ gives the stated
nef adjoint. The fiber formula follows from the two blowups; all other
fibers of the original ruling are unchanged and irreducible.
The graph projection to the $y$-line is $x\mapsto x^2$, and its strict
transform $B$ is isomorphic to the original graph. Thus this is an
inseparable bisection meeting $P_i$ once in a multiplicity-two
fiber. The other retained curves are vertical, so every horizontal exceptional curve has degree at most two.

For $n\ge4$, $B^2=4-2n<-2$, leaving exactly $F_i,U_i$ as isolated nodes.
In the integral total-transform basis,
\[
 F_i=b-E_{i1}-E_{i2},\qquad U_i=E_{i1}-E_{i2},\qquad
 F_i+U_i=b-2P_i.
\]
A subset $\sum_i(\alpha_iF_i+\beta_iU_i)$, with coefficients zero or one,
is twice a Picard class exactly when its coefficients in this free basis
are even. This is equivalent to
$\alpha_i=\beta_i$ for every $i$ and $\sum_i\alpha_i=0$ in $\FF_2$.
This identifies the code with $DE(n)$ in the coordinate order
$(F_1,U_1,\ldots,F_n,U_n)$ and proves the half-class formula; its
dimension and minimum weight are the standard ones stated above.
For $n=3$ include the coefficient $\gamma$ of
$B=2a+b-\sum_i(E_{i1}+E_{i2})$.
The parity equations are
$\beta_i=\alpha_i+\gamma$ and $\sum_i\alpha_i=\gamma$ in $\FF_2$.
They give dimension three. If $\gamma=0$, a nonzero word uses two pairs,
hence four nodes; if $\gamma=1$, it uses $B$ and exactly one node from
each of the three pairs, again four nodes.


\end{proof}

\begin{corollary}[Failure of the separability arguments in characteristic two]
\label{cor:inseparable-covers}
Both uses of odd characteristic in the upper-bound proof fail on the
preceding characteristic-two surfaces.
First, the bisection has degree-two fibers at every point and zero
differential, rather than two ramification points governed by a
separable Riemann--Hurwitz formula.
Second, for distinct indices $i,j$ there is the explicit even set
\[
 F_i+U_i+F_j+U_j\sim2(b-P_i-P_j).
\]
The degree-two algebra defined by its square-root equation is purely
inseparable, and its restriction to an original rational tree disjoint from
that set is a nonreduced double tree, not two disjoint copies.
Thus neither the bisection ramification bound nor the even-node
obstruction extends to characteristic two.
\end{corollary}

\begin{proof}
The bisection assertion follows from $y=x^2$ and the fiber
formula; by \cref{lem:bisection-ramification}, its
fiber multiplicities cannot be summed as separable ramification.
The even set was proved by the Picard-basis identity.
At its prime components the defining section has odd valuation, so it
is not a square in $k(S)$. Adjoining its square root in characteristic
two therefore gives a field extension that is purely inseparable of
degree two. On a disjoint original rational tree the half-class has
multidegree zero and is trivial by \cref{lem:tree-picard}; the section
restricts to a nonzero constant $c\in k$. The restricted algebra is
\[
 \mathcal O_T[z]/(z^2-c)
   \simeq\mathcal O_T[\varepsilon]/(\varepsilon^2),
\]
not $\mathcal O_T\oplus\mathcal O_T$.
For a function-field description, take $i,j$ to
be the selected points $x=0,1$. Away from the exceptional locus, the
square-root equation is $z^2=y(y-1)=y^2+y$.
Putting $w=z+y$ gives $w^2=y$, so the extension is purely
inseparable. Thus the doubled negative-definite subspace used
in the proof of \cref{thm:no-even-nodes} is absent.
\end{proof}


\begin{proof}[Proof of \cref{thm:examples}]
In \cref{prop:frobenius-family}, take $(p,n)=(3,3)$.
The contracted exceptional divisor has four isolated $(-3)$-curves
and three disjoint chains of two $(-2)$-curves. Thus it gives seven
distinct singular points, and the formula there gives $K_X^2=1/3$.
For $p=2$ and any $n\ge3$, the same proposition gives the family
$X_{2,n}$ and all the invariants stated in \cref{thm:examples}.
Together with \cref{thm:main}, the first example proves sharpness in
characteristic three.
\end{proof}

\appendix
\section{Negative curves on the resolution of the seven-point surface}
\label{app:negative-curves}
On the minimal resolution of the seven-point surface, the one-component
replacements at curves of least positive $L$-degree reduce the
singular-point count from seven to six.

Set $p=3$ and choose the three graph parameters
$I=\{0,1,\infty\}$, so the surface and all the following curves are defined
over $\FF_3$. Write $S=S_3$, $D=D_3$, and
\[
 L=a+\tfrac23b-\tfrac13\sum E_{ij}=\tfrac13(B+b),\qquad L^2=\tfrac13.
\]
When $\{i,j,k\}=I$, subscripts denote these three graph parameters. Every displayed equation denotes its
bihomogeneous closure and then its strict transform on $S$.

\begin{proposition}[The thirteen nonexceptional negative curves on the minimal resolution]
\label{prop:equality-negative-curves}
The complete list of exterior negative prime curves on $S$ is
\[
 \{P_i,T_i,Q_i,W_i:i\in I\}\ \cup\ \{\Theta\}.
\]
They are smooth rational $(-1)$-curves, with classes
\[
\begin{aligned}
 T_i&=a-E_{i1},& \Theta&=a+b-\sum_iE_{i1},\\
 Q_i&=2a+b-E_{i1}-\sum_{j\ne i}(E_{j1}+E_{j2}),\\
 W_i&=3a+2b-(E_{i1}+E_{i2}+E_{i3})
                    -\sum_{j\ne i}(2E_{j1}+E_{j2}).
\end{aligned}
\]
Here $T_i$ is the strict $a$-fiber through the selected point, and the
remaining equations are
\[
\begin{array}{c|l@{\qquad}c|l}
\Theta&x-y=0&&\\
Q_0&x^2+x+y=0&W_0&x^3+y^2+y=0\\
Q_1&x^2-y=0&W_1&x^3-y^2=0\\
Q_\infty&x^2+xy+y=0&W_\infty&x^3y+x^3+y^2=0.
\end{array}
\]
Their $L$-degrees are respectively
\[
 L\cdot P_i=L\cdot T_i=\tfrac13,\qquad
 L\cdot\Theta=L\cdot Q_i=\tfrac23,\qquad L\cdot W_i=1.
\]
In particular the six shortest exterior curves are exactly
$P_i,T_i$, and the least exterior degree is $1/3$, established without
using \cref{lem:nef-threshold} or the seven-point bound.
The surface has exactly twenty-three negative prime curves including $D$.
\end{proposition}

\begin{proof}
The charts in the preceding proposition construct $P_i,T_i$.
The diagonal and the three $Q_i$ are graphs of morphisms $\PP^1\to\PP^1$.
For $Q_1$, the graph $y=x^2$ has cluster multiplicities $(1,1,0)$ at
$0,\infty$ and $(1,0,0)$ at $1$.
The diagonal has $(1,0,0)$ at all three points.
The curve $W_1$ has normalization $x=z^2$, $y=z^3$ and only its two
cusps at $0,\infty$; each cusp is resolved by its first blowup.
Its cluster multiplicities are $(2,1,0)$ at those points and $(1,1,1)$
at $1$. For example at $0$, after writing $v=uv_1$ its strict equation
is $u=v_1^2$. Blowing up its common point with the strict graph, in the
other chart $u=v_1u_2$, gives $u_2=v_1$; this is smooth through the node
of $U_0$ and $V_0$. The third graph blowup is at a different point.
At $1$, the equations after the three graph substitutions are
$u^2+v_1-uv_1^2$, $u+v_2-u^2v_2^2$, and $1+v_3-u^3v_3^2$.
The last meets $P_1$ at $v_3=-1$, distinct from the three original contacts.
Reciprocal coordinates prove the same cusp calculation at infinity.
The simultaneous transformations $(x,y)\mapsto(g(x),g(y))$ for
$g(z)=1-z$ and $g(z)=z/(z-1)$ preserve the Frobenius graph and permute
$0,1,\infty$. They lift through the entire blowup configuration.
They carry the $Q_1,W_1$ calculations to the other two rows.
Thus every listed strict transform is smooth and rational.
The class formulas and intersection form give self-intersection $-1$
and the displayed $L$-degrees.

We prove completeness independently of those constructions.
Every exterior prime $Z$ has $K_S\cdot Z<0$; if $Z^2<0$, adjunction
makes it a smooth rational $(-1)$-curve.
Since $3L$ is integral and all four higher-weight coefficients are $1/3$,
write $L\cdot Z=r/3$ with $r=1,2,3$.
Let $z_0,z_i$ be its contacts with $B,F_i$, and $u_i,v_i$ its contacts
with $U_i,V_i$. The discrepancy-sum and projection identities give
\[
 z_0+\sum_i z_i=3-r,
 \qquad z_0^2+\sum_i z_i^2
       +2\sum_i(u_i^2+u_iv_i+v_i^2)=3+r^2.
\]
The negative matrix of $D$ is the sum of four $[3]$ blocks and three
$\left(\begin{smallmatrix}2&-1\\-1&2\end{smallmatrix}\right)$ blocks.
Solving its rank-one closure gives the coefficients of $[Z]$ in the
integral basis $a,b,E_{ij}$:
\[
\begin{aligned}
 [a]Z&=r-z_0,&[b]Z&=r-1,\\
 [E_{i1}]Z&=(-r+z_0+z_i-2u_i-v_i)/3,\\
 [E_{i2}]Z&=(-r+z_0+z_i+u_i-v_i)/3,\\
 [E_{i3}]Z&=(-r+z_0+z_i+u_i+2v_i)/3.
\end{aligned}
\]
In particular $z_0+z_i-r+u_i-v_i\equiv0\pmod3$ for each $i$.
For $r=1$, a single higher-weight contact two violates integrality at
an untouched branch. Otherwise two higher-weight contacts are one and
exactly one weight-two contact is one. At the two untouched branches
integrality forces $z_j=1-z_0$; at the touched branch it forces
$z_i=z_0$. For $z_0=0$ the weight-two contact is $U_i$, giving $T_i$;
for $z_0=1$ it is $V_i$, giving $P_i$.
For $r=2$, exactly one higher-weight contact is one and
$\sum(u_i^2+u_iv_i+v_i^2)=3$.
The congruences force all $u_i=1,v_i=0$ when $z_0=1$, giving $\Theta$;
when $z_i=1$ they force $u_i=1,v_i=0$ and $u_j=0,v_j=1$ for $j\ne i$,
giving $Q_i$.
For $r=3$ there are no higher-weight contacts. The quadratic equation gives
$u_i,v_i\le2$, and integrality forces $u_i=v_i$.
There must therefore be two pairs $(1,1)$ and one pair
$(0,0)$, giving exactly $W_i$ at the untouched branch.
All thirteen resulting integral classes were realized above.
Two distinct prime curves cannot share a negative numerical class,
since their nonnegative intersection would equal that negative square.
This proves completeness.
\end{proof}

\begin{proposition}[Adjoint identities and singular-point counts on the equality example]
\label{prop:equality-adjoints}
For $\{i,j,k\}=I$ the following are linear equivalences between actual
integral divisors on the unchanged $S$:
\[
\begin{aligned}
 K_S+V_i+B+F_i+2P_i&\sim T_i,\\
 K_S+U_i+F_j+F_k+2T_i&\sim P_i,\\
 K_S+\sum_iU_i+2\Theta&\sim\sum_iP_i,\\
 K_S+U_i+V_j+V_k+2Q_i&\sim P_i+T_j+T_k,\\
 K_S+U_j+V_j+W_i&\sim P_k+T_k\quad(j\ne i).
\end{aligned}
\]
The right sides are the unique effective representatives of their
classes. The six shortest curves thus form three degree-preserving
adjoint pairs $P_i\leftrightarrow T_i$; the other seven exterior
$(-1)$-curves have displayed adjoint reductions to shortest curves.
At $P_i$, the unique excess-contact vertex is $V_i$, of valency one in $D$;
at $T_i$ it is $U_i$, also of valency one in $D$.
The corresponding one-component replacement applies, but changes
the singular-point count by $1-2=-1$. Its target has six singular points,
minimal-resolution exceptional forest $3A_2+A_1+2[3]$, and anticanonical square $2/3$.
Thus these six one-component replacements do not provide a singular-point preserving
Picard descent.
\end{proposition}

\begin{proof}
The complete shortest contact rows are
\[
\begin{array}{c|c|c}
\text{curve}&\text{original contacts, all once}&\text{excess-contact vertex}\\\hline
P_i&B,F_i,V_i&V_i\\
T_i&F_j,F_k,U_i&U_i.
\end{array}
\]
The blowup charts give these intersections.
Substitution of the proved curve classes and
$K_S=-2a-2b+\sum E_{ij}$ gives all five adjoint identities.
Each right side consists of pairwise disjoint $(-1)$-curves with
coefficient one. Intersection with each of them forces it to occur in
any effective representative, and the remaining effective divisor would
be linearly trivial, hence zero. This proves the uniqueness assertions.
For the last identity the curve $W_i$ passes through the exceptional node
$U_j\cap V_j$ as a third smooth branch; the adjacent-contact adjoint
allows exactly this coincidence, so no distinct-point hypothesis is used.

Delete $V_i$ at $P_i$ and blow down $P_i$. The original $B$ and $F_i$
become transverse $(-2)$-curves forming one $A_2$ block. The curve $U_i$
remains a separate $A_1$ block, the two other $A_2$ blocks are unchanged,
and the other two $F$-curves remain isolated of square $-3$.
All contacts on $P_i$ were distinct, so these are the images.
The original retained support $P_i+(D-V_i)$ is an SNC forest and all
remaining contacts are bounded. Thus the one-component replacement
applies with $h=\eta=2/3$, $c=1$; its nonminimal anticanonical pullback is
\[
 L_i^\dagger=-K_S+P_i-\tfrac13(F_j+F_k).
\]
Its square is $2/3$, and the retained minimal-resolution configuration
is the forest with six connected components just listed. The component
change is $s_{V_i}-2=-1$.
At $T_i$ delete $U_i$ and blow down $T_i$ instead. The curves $F_j,F_k$
become the new $A_2$, $V_i$ is the isolated $A_1$, and $B,F_i$ are the two
remaining $[3]$ blocks. The nonminimal pullback is
\[
 \widetilde L_i^\dagger=-K_S+T_i-\tfrac13(B+F_i),
 \qquad (\widetilde L_i^\dagger)^2=\tfrac23.
\]
This proves the same conclusions. In both cases the contraction theorem
inside \cref{thm:one-component-replacement} verifies projectivity, klt singularities,
Picard number one, and ample anticanonical divisor of the target.

\end{proof}

\begin{corollary}[Intersection lattices at equality]
\label{cor:equality-lattice}
The seven-point bound of \cref{thm:main} is optimal.
The sharp surface has $K_S^2=-1$, ten original exceptional components,
three original edges, $\det A_D=3^7$, $\ell=L^2=1/3$, and at every shortest
center
\[
 p^{\mathsf T}A_D^{-1}p=\tfrac43,\qquad
 |\det\langle D,P\rangle|=3^6,\qquad
 [\Pic(S):\langle D,P\rangle]=27.
\]
It has no isolated exceptional nodes. Thus its full Picard index is integral, and there is no isolated-node
summand to which \cref{lem:picard-index} could apply.
For $\beta=3$, the single-adjoint table allows at most two edges,
as required by the assumption of at least eight singular points;
the example has three.
\end{corollary}

\begin{proof}
The contraction and seven singular points are proved in
\cref{prop:frobenius-family}; together with the independently
proved upper bound this gives optimality. The original matrix has four $[3]$
blocks and three $A_2$ blocks, giving determinant $3^7$ and no isolated
$A_1$ component. The projection and full-index formulas give the displayed
values. Finally $10-3=7$ is the component count of the exceptional forest.
\end{proof}

\clearpage
\begin{thebibliography}{99}
\bibitem{Belousov}
Grigory Belousov, \emph{The maximal number of singular points on log del
Pezzo surfaces}, J. Math. Sci. Univ. Tokyo \textbf{16} (2009), no.~2,
231--238.

\bibitem{Bernasconi}
Fabio Bernasconi, \emph{Kawamata--Viehweg vanishing fails for log del Pezzo
surfaces in characteristic $3$}, J. Pure Appl. Algebra \textbf{225}
(2021), no.~11, Article~106727.
\href{https://doi.org/10.1016/j.jpaa.2021.106727}{doi:10.1016/j.jpaa.2021.106727}.

\bibitem{CalabriCilibertoMendesLopes}
Alberto Calabri, Ciro Ciliberto, and Margarida Mendes Lopes,
\emph{Even sets of four nodes on rational surfaces}, Math. Res. Lett.
\textbf{11} (2004), 799--808.
\href{https://doi.org/10.4310/MRL.2004.v11.n6.a7}{doi:10.4310/MRL.2004.v11.n6.a7}.

\bibitem{CataneseLi}
Fabrizio Catanese and Binru Li,
\emph{Enriques' classification in characteristic $p>0$: the $P_{12}$-theorem},
Nagoya Math. J. \textbf{235} (2019), 201--226.
\href{https://doi.org/10.1017/nmj.2018.8}{doi:10.1017/nmj.2018.8}.

\bibitem{DolgachevMendesLopesPardini}
Igor Dolgachev, Margarida Mendes Lopes, and Rita Pardini,
\emph{Rational surfaces with many nodes}, Compos. Math.
\textbf{132} (2002), no.~3, 349--363.
\href{https://doi.org/10.1023/A:1016540925011}{doi:10.1023/A:1016540925011}.

\bibitem{Keel}
Se\'an Keel, \emph{Basepoint freeness for nef and big line bundles in positive
characteristic}, Ann. of Math. (2) \textbf{149} (1999), no.~1, 253--286.
\href{https://doi.org/10.2307/121025}{doi:10.2307/121025}.

\bibitem{KeelMcKernan}
Se\'an Keel and James McKernan, \emph{Rational Curves on Quasi-Projective
Surfaces}, Mem. Amer. Math. Soc. \textbf{140} (1999), no.~669.
\href{https://doi.org/10.1090/memo/0669}{doi:10.1090/memo/0669}.

\bibitem{KollarMori}
J\'anos Koll\'ar and Shigefumi Mori, \emph{Birational Geometry of Algebraic
Varieties}, Cambridge Tracts in Mathematics, vol.~134, Cambridge
University Press, Cambridge, 1998.
\href{https://doi.org/10.1017/CBO9780511662560}{doi:10.1017/CBO9780511662560}.

\bibitem{Lacini}
Justin Lacini, \emph{On rank one log del Pezzo surfaces in characteristic
different from two and three}, Adv. Math. \textbf{442} (2024),
Article~109568.
\href{https://doi.org/10.1016/j.aim.2024.109568}{doi:10.1016/j.aim.2024.109568}.

\bibitem{LiuXie}
Jihao Liu and Lingyao Xie, \emph{Number of singular points on projective
surfaces}, Chin. Ann. Math. Ser. B \textbf{46} (2025), no.~5, 713--724.
\href{https://doi.org/10.1007/s11401-025-0035-y}{doi:10.1007/s11401-025-0035-y}.

\bibitem{NagaokaFive}
Masaru Nagaoka, \emph{Non-log liftable log del Pezzo surfaces of rank one
in characteristic five}, Nagoya Math. J. \textbf{259} (2025), 399--422.
\href{https://doi.org/10.1017/nmj.2024.37}{doi:10.1017/nmj.2024.37}.

\bibitem{NagaokaReview}
Masaru Nagaoka, \emph{A review of Lacini's classification of rank one log
del Pezzo surfaces in characteristic different from two and three},
preprint, \href{https://arxiv.org/abs/2512.12294v1}{arXiv:2512.12294v1},
2025.

\bibitem{PalkaPelkaI}
Karol Palka and Tomasz Pe\l ka, \emph{Classification of del Pezzo
surfaces of rank one. I. Height 1 and 2. II. Descendants with elliptic
boundaries}, preprint, 2024, revised 2025,
\href{https://arxiv.org/abs/2412.21174v2}{arXiv:2412.21174v2}.

\bibitem{PalkaPelkaIII}
Karol Palka and Tomasz Pe\l ka, \emph{Classification of del Pezzo
surfaces of rank one. III. Height 3}, preprint,
\href{https://arxiv.org/abs/2508.13609v1}{arXiv:2508.13609v1}, 2025.

\bibitem{Stacks}
The Stacks Project Authors, \emph{The Stacks Project},
Section~53.12, \href{https://stacks.math.columbia.edu/tag/0C1B}{Tag~0C1B}
(Riemann--Hurwitz), accessed 5 September 2026.

\bibitem{Tanaka}
Hiromu Tanaka, \emph{Minimal model program for excellent surfaces},
Ann. Inst. Fourier (Grenoble) \textbf{68} (2018), no.~1, 345--376.
\href{https://doi.org/10.5802/aif.3163}{doi:10.5802/aif.3163}.
\end{thebibliography}
\end{document}
